% Document DPM1. Includes the angular comparison formerly maintained as ANM1.
\documentclass[11pt,letterpaper,reqno]{amsart}
% BEGIN INLINE manuscript-style
\makeatletter
\NeedsTeXFormat{LaTeX2e}[2020/10/01]
% Use with \documentclass[11pt,letterpaper,reqno]{amsart}.
% This file controls presentation only. It supplies no mathematical content.

\RequirePackage[T1]{fontenc}
\RequirePackage{lmodern}
\RequirePackage{microtype}
\RequirePackage[letterpaper,margin=1in,headheight=15pt]{geometry}
\RequirePackage{amsmath,amsthm,amssymb,mathtools}
\RequirePackage{enumitem}
\RequirePackage{perpage}
\RequirePackage{xcolor}
\definecolor{manuscriptred}{rgb}{0.6,0.1,0.1}
\definecolor{manuscriptblue}{rgb}{0.1,0.3,0.5}
\definecolor{manuscriptmagenta}{HTML}{60445C}
\RequirePackage[breaklinks=true,colorlinks=true,
  linkcolor=manuscriptred,citecolor=manuscriptblue,urlcolor=manuscriptmagenta]{hyperref}
\RequirePackage{cleveref}
\hypersetup{pdfauthor={}}

% Retain the AMS paragraph, title, heading and running-head conventions.
% Equations retain one global Arabic sequence, as in the reference.
\setcounter{tocdepth}{1}
\renewcommand{\thefigure}{\Roman{figure}}
\renewcommand{\thetable}{\Roman{table}}
\MakePerPage{footnote}
\renewcommand{\thefootnote}{%
  \ifnum\value{footnote}>9\arabic{footnote}\else\fnsymbol{footnote}\fi}

% Upright statements. Main results are A, B, C; supporting statements
% share a section-local sequence. These end marks indicate typography,
% never proof completeness, validation, or a deferred-proof status.
\theoremstyle{definition}
\newtheorem{ms@theorem}{Theorem}
\renewcommand{\thems@theorem}{\Alph{ms@theorem}}
\newtheorem*{ms@theoremstar}{Theorem}
\newtheorem{ms@proposition}{Proposition}[section]
\newtheorem{ms@lemma}[ms@proposition]{Lemma}
\newtheorem{ms@definition}[ms@proposition]{Definition}
\newtheorem{ms@corollary}[ms@proposition]{Corollary}
\newtheorem*{ms@corollarystar}{Corollary}

\newcommand{\ms@majorstart}{%
  \pushQED{\qed}\renewcommand{\qedsymbol}{$\blacksquare\square$}}
\newcommand{\ms@minorstart}{%
  \pushQED{\qed}\renewcommand{\qedsymbol}{$\blacksquare$}}
\newenvironment{theorem}[1][]{\ms@majorstart\ms@theorem[#1]}
  {\popQED\endms@theorem}
\newenvironment{theorem*}[1][]{\ms@majorstart\ms@theoremstar[#1]}
  {\popQED\endms@theoremstar}
\newenvironment{proposition}[1][]{\ms@minorstart\ms@proposition[#1]}
  {\popQED\endms@proposition}
\newenvironment{lemma}[1][]{\ms@minorstart\ms@lemma[#1]}
  {\popQED\endms@lemma}
\newenvironment{definition}[1][]{\ms@minorstart\ms@definition[#1]}
  {\popQED\endms@definition}
\newenvironment{corollary}[1][]{\ms@majorstart\ms@corollary[#1]}
  {\popQED\endms@corollary}
\newenvironment{corollary*}[1][]{\ms@majorstart\ms@corollarystar[#1]}
  {\popQED\endms@corollarystar}
% The source's lemma/proposition variants are deliberately not carried over:
% choose one marker for each statement type rather than encode proof status.

\theoremstyle{remark}
\newtheorem{remark}[ms@proposition]{Remark}
\newtheorem*{remark*}{Remark}
% The ordinary amsthm proof environment retains its open square.

\crefname{ms@theorem}{theorem}{theorems}
\Crefname{ms@theorem}{Theorem}{Theorems}
\crefname{ms@proposition}{proposition}{propositions}
\Crefname{ms@proposition}{Proposition}{Propositions}
\crefname{ms@lemma}{lemma}{lemmas}
\Crefname{ms@lemma}{Lemma}{Lemmas}
\crefname{ms@definition}{definition}{definitions}
\Crefname{ms@definition}{Definition}{Definitions}
\crefname{ms@corollary}{corollary}{corollaries}
\Crefname{ms@corollary}{Corollary}{Corollaries}
\crefname{remark}{remark}{remarks}
\Crefname{remark}{Remark}{Remarks}

% Bibliography configuration for this manuscript.
\makeatother
% END INLINE manuscript-style
\newcommand{\R}{\mathbb R}\newcommand{\C}{\mathbb C}
\newcommand{\Z}{\mathbb Z}\newcommand{\B}{\mathbb B}
\newcommand{\Ocub}{O_{\mathrm{cub}}}
\DeclareMathOperator{\supp}{supp}
\title[Diagonal positivity and rotation invariance]{Diagonal reflection positivity and\break rotation invariance without closedness}
\date{}
\hypersetup{hypertexnames=false,pdftitle={Diagonal reflection positivity and rotation invariance without closedness},pdfauthor={}}
\begin{document}
\begin{abstract}
A Gaussian-dominated law of distributional two-forms on $\R^3$, invariant
under translations and the cubic group and reflection positive at the
coordinate planes, is Euclidean invariant if it is reflection positive
at one diagonal plane. All its $L^2$ moments are covariant, including
at overlapping supports. We also prove a two-point result under weaker
symmetry assumptions, compare the angular criteria for scalar and
field-strength kernels, and give examples separating diagonal positivity
from the spectral cone and closedness.
\end{abstract}
\maketitle
\tableofcontents

\section{Introduction}\label{sec:introduction}

Let $\mu$ be a probability law on the real tempered two-forms
$\mathcal S'(\R^3;\Lambda^2)$, with the sigma-algebra generated by
evaluation on real Schwartz tests. Let $F$ be the canonical random
distribution and put
\[
 *f=(f_{23},-f_{13},f_{12}),\qquad \B(b)=F(*^{-1}b).
\]
For an isometry $\mathfrak g(x)=gx+a$ we use
\[
 (\mathfrak g\cdot b)(x)=\det(g)g\,b(\mathfrak g^{-1}x),\qquad
 (\mathfrak gF)(f)=F(\mathfrak g^{-1}\cdot f).
\]
Here the action on two-form tests is
$(\mathfrak g\cdot f)(x)=(\wedge^2g)f(\mathfrak g^{-1}x)$.
Reflection positivity has the characteristic-functional meaning in
\eqref{eq:rp} below, on tests compactly supported in the open positive
half space.

\begin{theorem}\label{thm:main}
Suppose that $\mu$ has the following properties:
\begin{enumerate}[label=\textup{(\roman*)}]
\item Gaussian domination with constant one:
\begin{equation}\tag{G}\label{eq:G}
 \mu(e^{\B(b)})\le e^{\|b\|_2^2/2}
 \quad\text{for every real Schwartz vector test }b;
\end{equation}
\item invariance under all translations and under the signed permutation
group $\Ocub=O(3)\cap\mathrm{GL}(3,\Z)$;
\item reflection positivity at $\{x_j=0\}$ for $j=1,2,3$, and at one
plane $\{x_i=\sigma x_l\}$, where $i<l$ and $\sigma\in\{1,-1\}$,
from either chosen side.
\end{enumerate}
Then $\mathfrak g_*\mu=\mu$ for every Euclidean isometry $\mathfrak g$.
For every $N\ge1$ the moment
\[
 S_N(b_1,\ldots,b_N)=\mu\bigl(\B(b_1)\cdots\B(b_N)\bigr)
\]
extends to a bounded complex multilinear form on
$L^2(\R^3;\C^3)^N$ and satisfies
\[
 S_N(\mathfrak g\cdot b_1,\ldots,\mathfrak g\cdot b_N)
   =S_N(b_1,\ldots,b_N)
\]
for all such tests, without a support restriction.
\end{theorem}

Hypothesis (H) means Gaussian domination, translation and cubic
invariance, and reflection positivity at the three coordinate planes.
The additional hypothesis of the theorem is positivity
at one diagonal plane. No condition $dF=0$ is imposed.

There are two obstacles to deriving the theorem from the spectral cone.
At two points, the angular growth estimate permits one extra frequency
for each of the equal circular labels. The law turned through $\pi/4$
has the same analytic bounds, while this extra coefficient has the
opposite sign in its reflection form. Positivity in both frames removes
the coefficient. The higher-order harmonic argument then proves
covariance at distinct planar points.

The second obstacle is contact. The moments exist on all $L^2$ tests,
so covariance away from coincidences does not yet prove the theorem.
We localize two entries in normalized discs of radius $d$. Gaussian
domination bounds their averaged products uniformly as $d\downarrow0$.
Weak limits of the associated pair fields define bounded fields over
the plane. The cone makes their
two-point kernels holomorphic, and the quarter turn joins the four
angular charts into a bounded entire function. Liouville's theorem
forces their reconstructed vectors into the zero-energy subspace.
In an ergodic component that subspace is the vacuum line. A further
covariance estimate is needed to show that the fields themselves are
constant in $L^2(\mu)$; equality of their reconstructed classes would
not suffice. This gives factorization of a contact pair. Induction then
removes the disjointness conditions one pair at a time.

The analytic tools used before this contact argument are the cone,
plate-kernel and harmonic arguments in \cite[Sections~7--8]{QEDCurrent}.
We recall the exact
forms needed below. Their hypotheses are conditions on the law; the
proof here does not require a lattice action or a choice of
renormalization parameters.
The proof of \cref{thm:main} is completed in \cref{sec:components},
where the covariance proved on the ergodic components is integrated
and the other coordinate planes are treated by cubic symmetry.

\subsection{Reflection forms, moments and the cone}

For a unit vector $n$ and $c\in\R$, let
$r_{n,c}x=x-2(n\cdot x-c)n$. Reflection positivity means
\begin{equation}\label{eq:rp}
 \sum_{a,b=1}^k\overline{z_a}z_b\,
 \mu\!\left(e^{i\{F(f_b)-F(r_{n,c}\cdot f_a)\}}\right)\ge0
\end{equation}
for every $k\ge1$, $z\in\C^k$, and real smooth two-form tests $f_a$
with compact support in $\{n\cdot x>c\}$. We denote its left side by
$\mathcal Q_\mu(n,c;z;f)$.
For a law invariant under $r_{n,c}$, the two choices of normal give
equivalent conditions. Reflection invariance is stated separately
whenever it is needed.

Write $\theta_j=r_{e_j,0}$. Conditions $(T)$, $(R_j)$ and $(RP_j)$
will mean translation invariance, invariance under $\theta_j$, and
positivity at $x_j=0$, respectively. For rotations set
\[
 R^{ji}_te_j=\cos t\,e_j+\sin t\,e_i,\qquad
 R^{ji}_te_i=-\sin t\,e_j+\cos t\,e_i.
\]
The third basis vector is fixed. Condition $(Q_{ji})$ means invariance
under $R^{ji}_{\pi/2}$. For a law with $(T)$, $(R_j)$ and $(RP_j)$,
the reflection-positive reconstruction has a class map
$[\,\cdot\,]_j:L^2(\Sigma_j^+,\mu)\longrightarrow\mathcal H_j$
with scalar product
\begin{equation}\label{eq:reconstructed-product}
 \langle[\Phi]_j,[\Psi]_j\rangle=\mu((\Theta_j\Phi)\Psi).
\end{equation}
Here $\Sigma_j^+$ is the sigma-algebra generated by evaluations on real
smooth tests compactly supported in $\{x_j>0\}$, and
$\Theta_j\Phi=\overline{\hat\theta_j\Phi}$, with
$(\hat{\mathfrak g}\Phi)(F)=\Phi(\mathfrak g^{-1}F)$.
The map is a contraction and has dense range; $\Omega_j=[1]_j$.
The nonnegative time generator $H_j$ and the commuting transverse
momenta have joint spectral measure $E_j$.
We write $(C_{ji})$ for the support condition
\begin{equation}\label{eq:cone}
 E_j\{(\omega,p_\perp):\omega<|p_i|\}=0,
 \quad\text{equivalently }H_j\ge|P_i|.
\end{equation}

The cone theorem \cite[Proposition~7.2]{QEDCurrent} says that
$(T),(R_j),(RP_j)$, together with reflection invariance and positivity
at the plane with normal $(e_j+e_i)/\sqrt2$, imply $(C_{ji})$.
It assumes no closedness. In the $(1,2)$-plane we abbreviate
\begin{equation}\tag{S}\label{eq:S}
 (G),\ (T),\ (R_1),\ (RP_1),\ (Q_{12}),\ (C_{12}).
\end{equation}

For $1\le p<\infty$ and integers $N\ge1$, Gaussian domination implies
the bounds
\begin{equation}\label{eq:moments}
 \|\B(b)\|_{L^p(\mu)}\le c_p\|b\|_2,\qquad
 |S_N(b_1,\ldots,b_N)|\le c_N^N\prod_a\|b_a\|_2,
 \qquad S_1=0,
\end{equation}
and $c_2=1$. These statements include the unique continuous extension
to $L^2$ tests. In particular,
\begin{equation}\label{eq:pair-bound}
 \|\B(b)\B(b')\|_{L^2(\mu)}\le c_4^2\|b\|_2\|b'\|_2.
\end{equation}
These bounds and the entire finite-dimensional generating functions
follow from \cite[Proposition~5.1]{QEDCurrent}.
The Taylor coefficients determine those generating functions, hence
the characteristic functions of all finite families of evaluations.
Thus invariance of every moment on real Schwartz tests is equivalent
to invariance of the law. Under a symmetry of the
law, the action on $\B$ extends to $L^2$ tests, and moment invariance
extends by \eqref{eq:moments}.

Let $\mathcal C_9$ denote the translation-invariant laws that are
invariant and reflection positive at every plane in the nine families
$x_j=c$ and $x_i\pm x_l=c$, $i<l$, $c\in\R$.

\begin{lemma}[Transport of diagonal positivity]\label{lem:transport}
Under (H), positivity at one diagonal plane through the origin is
equivalent to membership in $\mathcal C_9$.
\end{lemma}

\begin{proof}
Each coordinate or diagonal reflection through the origin is a signed
permutation. Hence (H) gives invariance under all of them. The signed
permutation group acts transitively on the twelve unit vectors
$(\pm e_i\pm e_l)/\sqrt2$, $i<l$. Conjugation gives
\[
 g r_{n,0}g^{-1}=r_{gn,0}.
\]
Changing variables in \eqref{eq:rp} therefore transports positivity
from the chosen diagonal normal to every such normal. Coordinate
positivity is already in (H). Finally,
\[
 r_{n,c}=\tau_{cn}r_{n,0}\tau_{-cn}
\]
and translation invariance transport both positivity and reflection
invariance to every parallel plane. This is precisely $\mathcal C_9$.
The converse follows from its definition.
\end{proof}

\section{Two-point covariance with one diagonal mirror}\label{sec:two-point}

The two-point assertion needs less symmetry than \cref{thm:main}.

Write $x=(y,x_3)$, $y=(y^0,y^1)=(x_1,x_2)$, $R_t=R^{12}_t$ and
$n_\upsilon=(\cos\upsilon,\sin\upsilon)$, also read as a vector in
$\R^3$. Set
\[
 e^\pm=e_1\pm ie_2,\quad e^3=e_3,\qquad
 w(\pm)=\pm1,\quad w(3)=0,\qquad s(\pm)=1,\quad s(3)=-1.
\]
A label $\lambda=(\sigma,\chi)$, with real $\chi\in L^2(\R)$, has sign
$s(\lambda)=s(\sigma)$ and, for $\alpha\in L^2(\R^2;\C)$, specifies the plate
$(\alpha\otimes\lambda)(x)=\alpha(y)\chi(x_3)e^\sigma$.

\begin{theorem}\label{thm:weak-two-point}
Suppose $\mu$ satisfies $(G),(T),(R_1),(RP_1),(Q_{12})$ and is
reflection positive at $x_1=x_2$ or at $x_1=-x_2$, from one chosen
side. Then
\[
 S_2(R^{12}_t\cdot b,R^{12}_t\cdot b')=S_2(b,b')
 \quad(t\in\R,\ b,b'\in L^2(\R^3;\C^3)).
\]
For labels $\lambda=(\sigma,\chi)$ and
$\lambda'=(\sigma',\chi')$ with real $\chi,\chi'\in C_c^\infty(\R)$,
the two-point plate kernel has the form
\[
 G_{\lambda\lambda'}(rn_\vartheta)
   =a_{\lambda\lambda'}(r)e^{iw\vartheta},\qquad
 |a_{\lambda\lambda'}(r)|\le\frac{72}{\pi}
       \|\chi\|_2\|\chi'\|_2 r^{-2},
\]
for every $r>0$ and $\vartheta\in\R$, where
$w=w(\sigma)+w(\sigma')$. The kernel represents $S_2$ on compact,
disjoint planar supports as in \eqref{eq:global-density}.
\end{theorem}

\subsection{The angular criterion}

For scalar kernels, rotation invariance means independence of the angle.
For two-form kernels, the vector labels also rotate, and covariance
requires a prescribed phase. The following criterion separates this
finite-dimensional calculation from the construction of the kernels.
Its second part will apply after the kernels have been joined.

\begin{proposition}[Elimination of the remaining angular modes]
\label{thm:criterion}
The following assertions hold.
\begin{enumerate}[label=\textup{(\roman*)}]
\item Let $h$ be entire and suppose, for some finite $A\ge0$, that
\[
 h(\zeta+\pi/2)=h(\zeta),\qquad
 |h(t+i\beta)|\le A e^{4|\beta|}
 \quad(t,\beta\in\R).
\]
If $h(i\beta)\ge0$ and $h(\pi/4+i\beta)\ge0$ for every
$\beta\in\R$, then $h$ is a nonnegative constant.
\item Let $w\in\{-2,-1,0,1,2\}$ and let $h$ be entire. Suppose, for
some finite $A\ge0$, that
\[
 h(\zeta+\pi/2)=e^{iw\pi/2}h(\zeta),\qquad
 |h(t+i\beta)|\le A e^{2|\beta|}
 \quad(t,\beta\in\R).
\]
For $w=\pm2$, assume in addition that
\[
 h(i\beta)\ge0,\qquad e^{-iw\pi/4}h(\pi/4+i\beta)\ge0
 \quad(\beta\in\R).
\]
Then $h(\zeta)=a e^{iw\zeta}$ for some $a\in\C$. For $w=\pm2$,
one has $a\ge0$.
\end{enumerate}
Each inequality includes the assertion that its left side is real.
\end{proposition}

The following Fourier coefficient argument is the one used in
\cite[Section~8.5]{QEDCurrent}, with an arbitrary growth exponent $q$.

\begin{lemma}\label{lem:frequencies}
Let $w\in\Z$ and $q\ge0$. Suppose that $h$ is entire and that,
for some finite $A\ge0$,
\[
 h(\zeta+\pi/2)=e^{iw\pi/2}h(\zeta),\qquad
 |h(t+i\beta)|\le A e^{q|\beta|}
 \quad(t,\beta\in\R).
\]
Then
\[
 h(\zeta)=
 \sum_{\substack{k\in\Z,\ |k|\le q\\ k-w\in4\Z}}a_k e^{ik\zeta},
 \qquad |a_k|\le A.
\]
\end{lemma}

\begin{proof}
The function $h$ is $2\pi$-periodic. For $k\in\Z$, its Fourier
coefficient can be written
\[
 a_k=\frac1{2\pi}\int_0^{2\pi}
             h(t+i\beta)e^{-ik(t+i\beta)}\,dt.
\]
The integral is independent of $\beta$: integrate the entire periodic
integrand around a rectangle and cancel the vertical sides. Hence
$|a_k|\le A e^{q|\beta|+k\beta}$. Letting $\beta\to-\infty$
for $k>q$ and $\beta\to+\infty$ for $k<-q$ removes those
coefficients. Translation by $\pi/2$ gives $a_k=0$ unless
$k-w\in4\Z$. Fourier uniqueness on the real axis and the identity
theorem give the finite expansion; $\beta=0$ gives its coefficient bound.
\end{proof}

\begin{proof}[Proof of \cref{thm:criterion}]
In the scalar case, \cref{lem:frequencies} leaves
$h(\zeta)=c_0+c_1e^{4i\zeta}+c_{-1}e^{-4i\zeta}$. Positivity gives
\[
 c_0+c_1e^{-4\beta}+c_{-1}e^{4\beta}\ge0,\qquad
 c_0-c_1e^{-4\beta}-c_{-1}e^{4\beta}\ge0.
\]
Divide by $e^{4\beta}$ and let $\beta\to+\infty$ to obtain
$c_{-1}\ge0$ and $-c_{-1}\ge0$. Thus $c_{-1}=0$.
Division by $e^{-4\beta}$ and passage to $\beta\to-\infty$
give $c_1=0$. The remaining inequality gives $c_0\ge0$.

In the tensor case, the conditions $|k|\le2$ and $k-w\in4\Z$
leave only $k=w$, except when $w=\pm2$. In those cases write
$h(\zeta)=a^w e^{iw\zeta}+a^{-w}e^{-iw\zeta}$.
Since $e^{-2iw\pi/4}=-1$, positivity reads
\[
 a^w e^{-w\beta}+a^{-w}e^{w\beta}\ge0,\qquad
 a^w e^{-w\beta}-a^{-w}e^{w\beta}\ge0.
\]
Multiply by $e^{-w\beta}$ and let $w\beta\to+\infty$.
The two limits give $a^{-w}\ge0$ and $-a^{-w}\ge0$.
Thus $a^{-w}=0$, and the first inequality then gives $a^w\ge0$.
\end{proof}

For a scalar plate kernel, such as a kernel associated with
$\operatorname{tr}F\cdot F$, the first part of \cref{thm:criterion}
applies once the following analytic inputs have been established:
the short-distance bound
$t^8\Gamma(t)\to0$, where $\Gamma$ is the Laplace transform of the
energy marginal of a positive joint energy--momentum measure; a
contour argument leaving only the frequencies $-4,0,4$; domination
$\widetilde\rho_{\mathrm d}\le C\widetilde\rho$ of the diagonal
joint measure by a measure supported on the cone; and identification
of the continued transforms with $h(i\beta)$ and
$h(\pi/4+i\beta)$. Given these inputs, the first part of
\cref{thm:criterion} removes the two nonconstant coefficients.
The two-form argument below instead obtains its growth exponent two
from Gaussian domination and retains the phase dictated by the vector
labels.

\subsection{Plates and the eight oriented mirrors}

Finite sums of plates are dense in $L^2(\R^3;\C^3)$, also with a
prescribed measurable support in the plane. Their norms and actions are
\begin{align}
 \|\alpha\otimes\lambda\|_2&\le\sqrt2\|\alpha\|_2\|\chi\|_2,
 \label{eq:plate-norm}\\
 R_t\cdot(\alpha\otimes\lambda)&=
 e^{-iw(\sigma)t}(\alpha\circ R_{-t})\otimes\lambda,\label{eq:plate-rotation}\\
 \Theta_1\B_\lambda(\alpha)&=
 s(\sigma)\B_\lambda(\bar\alpha\circ\theta_1).
 \label{eq:plate-reflection}
\end{align}
Here $\B_\lambda(\alpha)=\B(\alpha\otimes\lambda)$.
The plate identities and their use in the point and ordered-kernel
constructions are given in \cite[Sections~8.2--8.3]{QEDCurrent}.
The density assertion follows by approximating each component by
finite sums of products of a planar $L^2$ function and a transverse
$L^2$ function; the planar factors can be restricted to the prescribed
measurable set.

The group generated by $\theta_1$ and $R_{\pi/2}$ contains
$r_{n_{\pi/4},0}=R_{\pi/2}\theta_1$ and all its quarter-turn
conjugates. Thus the law is invariant under the reflections with
normals $n_\upsilon$, $\upsilon\in(\pi/4)\Z$.
Quarter-turn transport of coordinate positivity gives the normals
$\upsilon\in(\pi/2)\Z$; transport of the assumed diagonal positivity
gives the intervening normals. Both sides of each of the four planes
are included.

For each frame $\upsilon\in(\pi/4)\Z$ put
\[
 \mu_\upsilon=(R_{-\upsilon})_*\mu.
\]
Only two laws occur, since $\mu_{\upsilon+\pi/2}=\mu_\upsilon$.
In the frame $\upsilon$, the adjacent normals
$n_\upsilon,n_{\upsilon+\pi/4}$ supply coordinate and diagonal
reflection positivity. Orthogonal transport preserves $(G)$ and $(T)$
and conjugates the symmetry assumptions. The cone theorem therefore
gives $(C_{12})$ for every $\mu_\upsilon$. Its two-point function obeys
\begin{equation}\label{eq:turned-moment}
 S_2^\upsilon(b,b')=S_2(R_\upsilon\cdot b,R_\upsilon\cdot b').
\end{equation}

\subsection{Joining the kernels}

Throughout this kernel construction take real
$\chi,\chi'\in C_c^\infty(\R)$. General $L^2$ tests enter after the
kernel identity has been proved, through plate density and the moment
bound.

The one-picture plate-kernel construction in
\cite[Section~8.4]{QEDCurrent}
applies to each $\mu_\upsilon$: its assumptions are $(G),(T),(R_1),
(RP_1)$ and the cone on the one-field subspace, which follows from
$(C_{12})$. It gives a smooth kernel $G^\upsilon_{\lambda\lambda'}$
on $\{y^0>0\}$ and a holomorphic function
$h^\upsilon=h^{\upsilon,r}_{\lambda\lambda'}$ on
$|\Re\zeta|<\pi/2$. The statements we use are
\begin{align}
 S_2^\upsilon(\alpha\otimes\lambda,\alpha'\otimes\lambda')
 &=\iint\alpha(y)\alpha'(y')G^\upsilon_{\lambda\lambda'}(y'-y)\,dy\,dy',
 \label{eq:ordered-density}\\
 h^\upsilon(\vartheta)&=G^\upsilon_{\lambda\lambda'}(rn_\vartheta),
 \label{eq:angular-chart}\\
 |h^\upsilon(t+i\beta)|&\le\frac{36}{\pi}\|\chi\|_2\|\chi'\|_2
       (r\cos t)^{-2}e^{2|\beta|},\label{eq:angular-bound}\\
 s(\sigma)h^\upsilon_{\lambda\lambda}(i\beta)&\ge0.
 \label{eq:angular-positive}
\end{align}
The first formula requires compact supports ordered by
$\sup_{\supp\alpha}y^0<\inf_{\supp\alpha'}y^0$; the next two use
$|\vartheta|,|t|<\pi/2$. The last holds for every real $\beta$.

Fix the two labels and put $w=w(\sigma)+w(\sigma')$.
The kernel in the original coordinates of frame $\upsilon$ is
\[
 G^{(\upsilon)}(y)=e^{iw\upsilon}
     G^\upsilon_{\lambda\lambda'}(R_{-\upsilon}y),
 \qquad n_\upsilon\cdot y>0.
\]
By \eqref{eq:plate-rotation} and \eqref{eq:turned-moment}, it
represents the original $S_2$ when the supports are ordered along
$n_\upsilon$. If two frames apply, the resulting integrals agree for
all tests in sufficiently small compact neighborhoods of any pair
in the common domain. Equality of continuous densities gives
$G^{(\upsilon)}=G^{(\upsilon')}$ there. These half planes cover
$\R^2\setminus\{0\}$, so the kernels define a single smooth
$G_{\lambda\lambda'}$ on that set.

The density formula holds for every pair of compact, disjoint supports:
\begin{equation}\label{eq:global-density}
 S_2(\alpha\otimes\lambda,\alpha'\otimes\lambda')
 =\iint\alpha(y)\alpha'(y')G_{\lambda\lambda'}(y'-y)\,dy\,dy'.
\end{equation}
To see this, partition both supports into sufficiently small squares.
Their positive distance allows each pair of squares to be strictly
ordered in one of the four coordinate directions. Apply the ordered
formula to each pair and sum; only finitely many squares meet the
supports. All integrals converge absolutely because the supports are
compact, the tests are in $L^1$ there, and the kernel is bounded on
their set of differences. The density is unique. It is consequently
bilinear in the real profiles and satisfies
\begin{equation}\label{eq:kernel-symmetry}
 G_{\lambda\lambda'}(y)=G_{\lambda'\lambda}(-y).
\end{equation}

For a fixed $r>0$, define on $|\Re\zeta-\upsilon|<\pi/2$
\[
 h^{(\upsilon)}(\zeta)=e^{iw\upsilon}h^\upsilon(\zeta-\upsilon).
\]
On real overlaps these functions equal
$G_{\lambda\lambda'}(rn_\vartheta)$. The identity theorem joins them
to an entire function $h$. Choosing a frame with
$|t-\upsilon|\le\pi/4$ in \eqref{eq:angular-bound} gives
\begin{equation}\label{eq:entire-bound}
 |h(t+i\beta)|\le\frac{72}{\pi}\|\chi\|_2\|\chi'\|_2
       r^{-2}e^{2|\beta|}.
\end{equation}
Since the chart kernels repeat after $\pi/2$, one also has
$h(\zeta+\pi/2)=e^{iw\pi/2}h(\zeta)$. \Cref{lem:frequencies} now gives
\begin{equation}\label{eq:two-frequencies}
 h(\zeta)=\sum_{\substack{|k|\le2\\ k-w\in4\Z}}a^k(r)e^{ik\zeta}.
\end{equation}
Each coefficient is bounded by the right side of
\eqref{eq:entire-bound} with $\beta=0$.

\subsection{The sign at the diagonal frame}

Only the equal circular labels $\sigma=\sigma'=\pm$ admit a
frequency other than $w$: then $w=\pm2$ and the other frequency
is $-w$. Start with equal profiles, so $\lambda'=\lambda$.
Since $s(\sigma)=1$, \eqref{eq:angular-positive} in frame $\upsilon$
reads
\[
 e^{-iw\upsilon}h(\upsilon+i\beta)\ge0.
\]
The frames $0$ and $\pi/4$ give the two positivity assumptions in
\cref{thm:criterion}. Together with \eqref{eq:entire-bound} and the
quarter-turn relation, its second part gives $a^{-w}=0$.
For unequal profiles the coefficient is bilinear and, by
\eqref{eq:kernel-symmetry}, symmetric at the even frequency $-w$.
Polarization gives its vanishing for those profiles as well.

It follows that $h(\zeta)=a^w(r)e^{iw\zeta}$ for every pair of labels,
and thus
\begin{equation}\label{eq:two-covariance}
 G_{\lambda\lambda'}(R_ty)=e^{iwt}G_{\lambda\lambda'}(y).
\end{equation}
This proves the kernel assertion of \cref{thm:weak-two-point}.
Substitution in \eqref{eq:global-density} cancels the phase from
\eqref{eq:plate-rotation}. Hence the rotation defect
\[
 D_t(b,b')=S_2(R_t\cdot b,R_t\cdot b')-S_2(b,b')
\]
vanishes for tests over disjoint compact planar sets, first for plates
and then by density. It remains bounded on the whole space:
\begin{equation}\label{eq:two-defect-bound}
 |D_t(b,b')|\le2\|b\|_2\|b'\|_2.
\end{equation}

\subsection{The two-point contact defect}

We need the following elementary form of locality.

\begin{lemma}\label{lem:local-constant}
Let $E$ be a bounded bilinear form on $L^2(\R^2;\C)$. If it is
translation invariant and vanishes on tests supported in disjoint
compact sets, then $E(\alpha,\alpha')=c\int\alpha\alpha'$ for a
constant $c\in\C$.
\end{lemma}

\begin{proof}
For the dyadic squares $Q_{n,k}=2^{-n}(k+[0,1)^2)$, the form vanishes
on distinct squares. Indeed, shrink each square to a compact interior
square, apply the hypothesis, and pass to the limit in $L^2$.
Translation invariance makes $E(1_{Q_{n,k}},1_{Q_{n,k}})$ independent
of $k$. Splitting a square into four children shows that this value
is $4^{-n}c$, where $c=E(1_{[0,1)^2},1_{[0,1)^2})$.
The required formula follows on finite dyadic step functions.
Such functions are dense in $L^2$, and both sides are bounded
bilinear forms, so it holds everywhere.
\end{proof}

Apply the lemma, for fixed labels and $t$, to
$E_t(\alpha,\alpha')=D_t(\alpha\otimes\lambda,
\alpha'\otimes\lambda')$. Translation invariance follows from
$R_t\tau_v=\tau_{R_tv}R_t$ and $(T)$. Thus
\[
 E_t(\alpha,\alpha')=c_t(\lambda,\lambda')\int\alpha\alpha'.
\]
Under simultaneous rotation of the two plates this expression acquires
the phase $e^{-iws}$. Quarter-turn invariance of $S_2$ makes it
unchanged at $s=\pi/2$, so $c_t=0$ unless $w=0$. Since $|w|\le2$,
these are all the possibilities. In either case $D_t$ is invariant
under every simultaneous rotation, first on plates and then by density
on all tests.

The rotation defect satisfies the cocycle identity
\[
 D_{t+s}(b,b')=D_t(R_s\cdot b,R_s\cdot b')+D_s(b,b').
\]
Its just-proved invariance makes this $D_{t+s}=D_t+D_s$. For every
integer $k\ge1$, \eqref{eq:two-defect-bound} therefore gives
\[
 k|D_t(b,b')|=|D_{kt}(b,b')|\le2\|b\|_2\|b'\|_2.
\]
Letting $k\to\infty$ proves $D_t=0$ and completes the proof of
\cref{thm:weak-two-point}.

\subsection{Comparison with closedness}

There is another way to remove the same exceptional coefficient.
Suppose that a law satisfies \textup{(S)} and is closed, so that
$\B(\nabla u)=0$ for scalar Schwartz tests $u$. Under these assumptions,
the two-point kernel construction in
\cite[Section~8.5]{QEDCurrent} gives the same frequency window as
\eqref{eq:two-frequencies}.
For real $\chi_0,\chi'\in C_c^\infty(\R)$ and
$\lambda'=(\sigma',\chi')$, put
$z=y^0+iy^1$ and
$\partial_z=(\partial_{y^0}-i\partial_{y^1})/2$.
Take $\eta,\alpha'\in C_c^\infty(\R^2;\C)$ with disjoint supports.
Pairing
\[
 \nabla(\eta(y)\chi_0(x_3))
   =(\partial_z\eta)\otimes(+,\chi_0)
      +(\partial_{\bar z}\eta)\otimes(-,\chi_0)
      +\eta\otimes(3,\chi_0')
\]
with the plate $\alpha'\otimes\lambda'$, and integrating by
parts in the density formula, gives
\begin{equation}\label{eq:bianchi}
 \partial_zG_{(+,\chi_0)\lambda'}
 +\partial_{\bar z}G_{(-,\chi_0)\lambda'}
 +G_{(3,\chi_0')\lambda'}=0
 \quad\text{on }\R^2\setminus\{0\}.
\end{equation}
For a radial coefficient $f$,
\[
 \partial_z(fe^{ik\vartheta})
   =\tfrac12e^{i(k-1)\vartheta}(f'+kf/r),\qquad
 \partial_{\bar z}(fe^{ik\vartheta})
   =\tfrac12e^{i(k+1)\vartheta}(f'-kf/r).
\]
When $\sigma'=+$, only the exceptional coefficient
$f=a^{-2}_{++}$ contributes to frequency $-3$ in
\eqref{eq:bianchi}; when $\sigma'=-$, only $f=a^2_{--}$
contributes to frequency $3$. In both cases it follows that
$f'-2f/r=0$, hence $f(r)=Cr^2$. The coefficient bound
$|f(r)|\le c r^{-2}$ holds for every $r>0$, and passage to
$r\to\infty$ gives $C=0$. The use of the estimate at all radii
is necessary here.

Closedness and diagonal positivity therefore supply different
hypotheses for eliminating the exceptional frequency. Rotational
covariance itself does not imply closedness; \cref{prop:nonclosed}
gives an invariant law positive at every plane for which $dF\ne0$.
At coincidences, the bounded-form theorem of
\cite[Theorem~A]{Coincidences}
requires vanishing both on entries over pairwise disjoint compact
planar sets and whenever any entry is a compactly supported gradient.
The proof in \cref{sc:section} uses localized pairs to remove the
support restriction without this gradient condition.

\section{All moments away from coincidences}\label{sec:harmonics}

We next use the higher-order harmonic constructions and estimates of
\cite[Sections~8.7--8.9]{QEDCurrent}. Their proofs use exactly (S);
the closedness argument at two points is separate. For $N\ge2$ and a tuple of
labels $\lambda=(\lambda_1,\ldots,\lambda_N)$, they give real analytic
functions on
\[
 \mathrm{Conf}_N(\R^2)=\{(y_1,\ldots,y_N):y_a\ne y_b\ (a\ne b)\}
\]
with finite expansion
\begin{equation}\label{eq:all-harmonics}
 \widetilde G_N(R_tY;\lambda)
 =\sum_m e^{i(w+4m)t}\widetilde c_m(Y;\lambda),
 \qquad w=\sum_aw(\sigma_a).
\end{equation}
The distinct-configuration density argument in
\cite[Section~8.10]{QEDCurrent} identifies these
kernels as the moment densities on plates with compact, pairwise
disjoint planar supports. Its proof moves one of two time-separated
groups transversely. The cone makes the correlation analytic, and
the identity theorem extends the ordered density formula to the
configuration space of distinct planar points. This identification
is needed in addition to the angular expansion. It does not remove
terms supported at coincidences.

For $N\ge2$, the anomaly index $\mathfrak a(N)$ is the largest
$|m|$ for which $\widetilde c_m(\,\cdot\,;\lambda)$ is not identically
zero for some tuple of labels. It is zero if every harmonic vanishes;
set $\mathfrak a(0)=\mathfrak a(1)=0$. The harmonic and polygon
results give, for nonnegative integers $a,b$ and $N$,
\begin{align}
 \mathfrak a(0)&=0,\qquad \mathfrak a(N)\le N/2,\label{eq:anomaly-linear}\\
 2\mathfrak a(a+b)&\le\mathfrak a(2a)+\mathfrak a(2b),
 \label{eq:anomaly-convex}\\
 a&\mid 2\mathfrak a(2a)\quad(a\ge1).
 \label{eq:anomaly-polygon}
\end{align}
Complex conjugation exchanges the circular labels and replaces the
index $m$ by $-m$. Whenever $\mathfrak a(N)>0$, both extreme indices
occur. In particular, if
$\mathfrak a(2)=1$, some two-point harmonic $\widetilde c_{-1}$ is
nonzero. The top-harmonic lemma gives the positive
semidefinite kernel
\begin{equation}\label{eq:top-kernel}
 K_1((x;\lambda),(z;\lambda'))
 =s(\lambda)\widetilde c_{-1}(\theta_1x,z;\lambda,\lambda'),
 \qquad x,z\in\{y^0>0\}.
\end{equation}
The harmonic extension, convexity, regular-polygon argument and
positive top-coefficient kernel are proved in
\cite[Sections~8.7--8.9]{QEDCurrent}. These arguments do not use
$dF=0$.

\begin{proposition}\label{prop:separated}
Suppose $\mu$ satisfies $(G),(T),(R_1),(RP_1),(R_2),(RP_2),(Q_{12})$
and is invariant and reflection positive at $x_1+x_2=0$.
Then $\mathfrak a(N)=0$ for every $N$, and all its moments are
covariant under $R_t$ on tests over pairwise disjoint compact planar
sets.
\end{proposition}

\begin{proof}
The cone theorem gives (S) for $\mu$. Let $g=R_{\pi/4}$ and
$\mu'=(g^{-1})_*\mu$. The map $g$ sends the coordinate normal $e_1$
to $n_+=(e_1+e_2)/\sqrt2$ and sends $n_+$ to $e_2$.
Thus $\mu'$ has the coordinate and diagonal symmetries and positivity
needed for the cone theorem as well. It has $(G),(T),(Q_{12})$ by
transport, so it too satisfies (S).

The moment transformation formula and uniqueness of the plate
densities give
\[
 \widetilde G'_N(Y;\lambda)
   =e^{-iw\pi/4}\widetilde G_N(R_{\pi/4}Y;\lambda).
\]
Use \eqref{eq:all-harmonics} and compare the finite Fourier expansions
after replacing $Y$ by $R_tY$. This gives
\begin{equation}\label{eq:sign-change}
 \widetilde c'_m=(-1)^m\widetilde c_m.
\end{equation}
If $\mathfrak a(2)=1$, a nonzero $\widetilde c_{-1}$ remains nonzero
for $\mu'$, so its index is also one. The top-harmonic kernel for
$\mu'$ is then $-K_1$. Both $K_1$ and $-K_1$ are positive
semidefinite, which forces their diagonal and hence the entire kernel
to vanish. Thus $\widetilde c_{-1}$ vanishes on
$(-\mathbb H)\times\mathbb H$, where $\mathbb H=\{y^0>0\}$.
Real analyticity on the connected configuration space forces its
vanishing everywhere, a contradiction. The bound
$\mathfrak a(2)\le1$ therefore gives $\mathfrak a(2)=0$.

For $\nu(a)=\mathfrak a(2a)$, \eqref{eq:anomaly-linear}--
\eqref{eq:anomaly-polygon} give a nonnegative integer sequence with
$\nu(0)=0$, $\nu(a)\le a$, nondecreasing increments and
$a\mid2\nu(a)$. The integer-sequence argument in
\cite[Section~8.9]{QEDCurrent} therefore gives
$\mathfrak a(2a)=a\mathfrak a(2)$.
For arbitrary $N$, choose $a=\lfloor N/2\rfloor$ and $b=N-a$
in \eqref{eq:anomaly-convex}. Thus
\begin{equation}\label{eq:reduction}
 \mathfrak a(N)\le\frac N2\mathfrak a(2).
\end{equation}
The vanishing of the two-point index therefore gives
$\mathfrak a(N)=0$ for every $N$.

Consequently \eqref{eq:all-harmonics} reduces to
\[
 \widetilde G_N(R_tY;\lambda)=e^{iwt}\widetilde G_N(Y;\lambda).
\]
In the density formula this phase cancels the product of the plate
phases in \eqref{eq:plate-rotation}. This proves covariance for plates
over pairwise disjoint compact sets. Plate density with prescribed
supports and \eqref{eq:moments} give it for all tests over those sets.
\end{proof}

\Cref{prop:separated} proves the separated-support assertion
needed for \cref{thm:main}. We now remove the
support restriction.


\section{Contact pairs and invariance on all test fields}
\label{sc:section}

The angular argument gives covariance on entries over pairwise disjoint
compact subsets of the plane. To remove one disjointness condition,
we localize the corresponding pair in the same disc. Gaussian
domination gives pair fields bounded from $L^1(\R^2)$ to $L^2(\mu)$,
uniformly in the disc radius. A weak limit has a bounded reconstructed
two-point kernel on the forward tube. Quarter-turn covariance and
Liouville's theorem put its reconstructed classes in the zero-energy
space. When that space is the vacuum line, reflection factorization
and an $L^1$ strip estimate show that the field itself is constant in
$L^2(\mu)$. This yields factorization of the localized pair; induction
then removes the disjointness conditions one pair at a time.

Throughout this section write
\[
\mathfrak h=L^2(\mathbb R^3;\mathbb C^3),\qquad
x=(y,x_3),\quad y=(x_1,x_2),\qquad R_t=R^{12}_t,
\]
and let $R=R_{\pi/2}$. The assumptions (S) are those of \eqref{eq:S}. The unadorned norm on $\mathfrak h$ is its $L^2$ norm.
The additional condition $\ker H_1=\mathbb C\Omega_1$ will be stated
where it is used. Here $\mathcal H_1$, $[\,\cdot\,]_1$, $H_1$
and $\Omega_1=[1]_1$ are the reconstruction in the first coordinate.
For a Borel set $K\subset\mathbb R^2$ set
\[
\mathfrak h_K=\{f\in\mathfrak h:f=0\text{ a.e. outside }K\times\mathbb R\}.
\]
A bounded function $a$ on the plane acts on $\mathfrak h$ by multiplication
in $y$. In particular, $1_K$ is the orthogonal projection onto $\mathfrak h_K$.

We use the Gaussian bounds \eqref{eq:moments} and \eqref{eq:pair-bound};
H\"older's inequality also puts every finite product of fields in
$L^2(\mu)$. Write $U(v)=\hat\tau_v$ for planar translations.
Symmetries act unitarily and multiplicatively on $L^2(\mu)$ and
commute with conjugation; $\Theta_1$ is an antiunitary involution.
The reconstructed scalar product was defined in \cref{sec:introduction}, and
\begin{equation}\label{sc:reconstruction}
 \langle\Omega_1,[\Phi]_1\rangle=\mu(\Phi),\qquad
 \|[\Phi]_1\|\le\|\Phi\|_{L^2(\mu)}
 \quad(\Phi\in L^2(\Sigma_1^+,\mu)).
\end{equation}
All scalar products are antilinear in their first argument.

\subsection{Localization to one pair}

For $d>0$ and $z\in\mathbb R^2$ let
\[
\varphi_{d,z}(y)=(\pi d^2)^{-1/2}1_{B(z,d)}(y).
\]
The normalization and covariance identities are
\begin{align}
&\|\varphi_{d,z}\|_2=1,\qquad
 \int\varphi_{d,z}(y)^2\,\mathrm dz=1,\qquad
 \int\|\varphi_{d,z}f\|^2\,\mathrm dz=\|f\|^2,
 \label{sc:discs}\\
&\varphi_{d,z}\circ\theta_1=\varphi_{d,\theta_1z},\qquad
 \varphi_{d,z}\circ R_{-t}=\varphi_{d,R_tz},\qquad
 \tau_v\varphi_{d,z}=\varphi_{d,z+v},\qquad
 R_t\cdot(\varphi_{d,z}f)=\varphi_{d,R_tz}(R_t\cdot f).
 \label{sc:disc-actions}
\end{align}
The third identity in \eqref{sc:discs} follows from the second by Tonelli's
theorem. Both $z\mapsto\varphi_{d,z}$ in $L^2(\mathbb R^2)$ and
$z\mapsto\varphi_{d,z}f$ in $\mathfrak h$ are norm continuous by
dominated convergence, since the boundary circle has measure zero.

\begin{lemma}[Localization of a bilinear form]\label{sc:localization}
Let $K_1,K_2\subset\mathbb R^2$ be Borel, and let
$\beta:\mathfrak h_{K_1}\times\mathfrak h_{K_2}\to\mathbb C$
be bounded and bilinear. Suppose $\beta(f_1,f_2)=0$ whenever
$f_a\in\mathfrak h_{K_a\cap A_a}$ for two disjoint compact sets
$A_1,A_2\subset\mathbb R^2$. Then, for every Borel set $A$,
\begin{equation}\label{sc:locality}
\beta(1_Af_1,1_Af_2)=\beta(1_Af_1,f_2),
\end{equation}
and, for every $d>0$,
\begin{equation}\label{sc:localization-identity}
\beta(f_1,f_2)=\int\beta(\varphi_{d,z}f_1,\varphi_{d,z}f_2)\,\mathrm dz.
\end{equation}
The integral in \eqref{sc:localization-identity} is absolutely convergent.
\end{lemma}

\begin{proof}
The assumed vanishing holds also for two disjoint Borel sets $A,A'$.
To see this, exhaust each by increasing compact subsets up to a null set,
using inner regularity on its intersections with bounded discs. The
corresponding multiplications converge strongly on $\mathfrak h$, so
continuity of $\beta$ gives $\beta(1_Af_1,1_{A'}f_2)=0$.
Splitting $f_2=1_Af_2+1_{A^c}f_2$ proves \eqref{sc:locality}.

For fixed $f_2$, the Riesz representation theorem gives
$u\in\mathfrak h_{K_1}$ such that
$\beta(f,f_2)=\int\overline{u(x)}\cdot f(x)\,\mathrm dx$.
By \eqref{sc:locality},
\[
\beta(\varphi_{d,z}f_1,\varphi_{d,z}f_2)
 =\int\varphi_{d,z}(y)^2\overline{u(x)}\cdot f_1(x)\,\mathrm dx.
\]
Since $\overline u\cdot f_1$ is integrable, Fubini's theorem and
\eqref{sc:discs} prove both assertions about the integral.
\end{proof}

We use Hilbert-space integrals in their weak sense. If $V(z)$ is norm
continuous with values in a Hilbert space and
$\int|a(z)|\|V(z)\|\,\mathrm dz<\infty$, its integral is the unique vector
$v$ satisfying
\[
\langle\eta,v\rangle=\int a(z)\langle\eta,V(z)\rangle\,\mathrm dz.
\]
Existence follows from the Riesz theorem \cite[Theorem 3.3]{Melrose}, and
$\|v\|\le\int|a|\|V\|$. Bounded linear maps pass through this integral;
an antiunitary involution $J$ gives
$Jv=\int\overline a\,JV$. An integral of vectors in a closed subspace
belongs to that subspace, by orthogonality to its complement. In $L^2(\mu)$
we shall repeatedly use
\begin{equation}\label{sc:integral-pairing}
\mu\left(\Phi\int aV\right)=\int a(z)\mu(\Phi V(z))\,\mathrm dz
\qquad(\Phi\in L^2(\mu)).
\end{equation}

For complex angles we use the complex linear rotation $R_\zeta$ defined by
$(R_\zeta\xi)^\pm=e^{\pm i\zeta}\xi^\pm$, where
$\xi^\pm=\xi^0\pm i\xi^1$. Thus
$R_{i\beta}(\ell,0)=(\ell\cosh\beta,i\ell\sinh\beta)$.

\subsection{Bounded local fields have only vacuum classes}

The limiting pair will be described by a bounded map
$X:L^1(\mathbb R^2;\mathbb C)\to L^2(\mu)$. Write $L^1_+$ for the
functions vanishing a.e. on $\{y^0<0\}$ and $L^1_-$ for those vanishing
a.e. on $\{y^0>0\}$. We require, for some $C\ge0$,
$s\in\{1,-1\}$ and $w\in\mathbb Z$, that
\begin{align}
&X\text{ is complex linear},\qquad \|X(h)\|_{L^2(\mu)}\le C\|h\|_1,
 \label{sc:bounded-field}\\
&U(v)X(h)=X(\tau_vh)\quad(v\in\R^2,\ h\in L^1),
 \label{sc:field-translation}\\
&X(h)\in L^2(\Sigma_1^+,\mu)\quad(h\in L^1_+),
 \label{sc:field-locality}\\
&\Theta_1X(h)=sX(\overline h\circ\theta_1),\qquad
 \hat RX(h)=(-i)^wX(h\circ R^{-1})\quad(h\in L^1).
 \label{sc:field-actions}
\end{align}
Here $(\tau_vh)(y)=h(y-v)$.

\begin{lemma}[Vacuum classes]\label{sc:vacuum-classes}
Assume \textup{(S)}, and let $X$ satisfy
\eqref{sc:bounded-field}--\eqref{sc:field-actions}.
There is a vector $v_0\in\ker H_1$ such that
\[
[X(h)]_1=\left(\int h\right)v_0\qquad(h\in L^1_+).\qedhere
\]
\end{lemma}

\begin{proof}
Let $E_1$ be the joint spectral measure of the reconstruction and
put
\[
\mathcal W=\{\xi\in\mathbb C^2:\Re\xi^0>|\Im\xi^1|\},\qquad
\mathsf T(\xi)=\int e^{-\xi^0\omega+i\xi^1p_2}\,
 \mathrm dE_1(\omega,p_\perp).
\]
The cone condition means that $E_1$ is carried by
$\{\omega\ge|p_2|\}$. The cone semigroup is therefore a contraction on
$\overline{\mathcal W}$, strongly continuous there and holomorphic in
operator norm on $\mathcal W$. It satisfies
\[
\mathsf T(\xi)\mathsf T(\xi')=\mathsf T(\xi+\xi'),\qquad
\mathsf T(\xi)^*=\mathsf T(\xi^*),\quad
\xi^*=(\overline{\xi^0},-\overline{\xi^1}),
\]
and $\mathsf T(y)[\Phi]_1=[U(y)\Phi]_1$ for real $y$ with $y^0\ge0$.
These conclusions follow from the spectral-calculus argument in
\cite[Section~8.2]{QEDCurrent}: the multiplier has modulus at most
one on the closed tube, and on a smaller tube its first two
derivatives are bounded by constants times
$\omega e^{-\varepsilon\omega}$ and
$\omega^2e^{-\varepsilon\omega}$. Dominated convergence gives
strong continuity, and Taylor's formula gives operator-norm
holomorphy.

Set $\psi(h)=[X(h)]_1$ for $h\in L^1_+$. The contraction in
\eqref{sc:reconstruction} and the bound \eqref{sc:bounded-field}
give the first assertion below. The semigroup identity just stated,
together with \eqref{sc:field-translation}, gives the second:
\begin{equation}\label{sc:psi}
\|\psi(h)\|\le C\|h\|_1,\qquad
\mathsf T(y)\psi(h)=\psi(\tau_yh)\quad(y^0\ge0).
\end{equation}
Take $g_n=n^2 1_{(0,1/n)^2}$. The vectors $\psi(g_n)$ are bounded, so a
subsequence converges weakly to a vector $v_0$ with $\|v_0\|\le C$.
Apply diagonal extraction in the separable closed span of this bounded
sequence and then the Riesz theorem
\cite[Proposition 3.13 and Lemma 3.11]{Melrose}; no separability of the
whole Hilbert space is needed.

For $h\in L^1_+$ we claim
\begin{equation}\label{sc:point-vector}
\langle\eta,\psi(h)\rangle
 =\int h(y)\langle\eta,\mathsf T(y)v_0\rangle\,\mathrm dy.
\end{equation}
Indeed,
\[
\|h*g_n-h\|_1
 \le\sup_{|a|\le2/n}\|\tau_ah-h\|_1\longrightarrow0.
\]
For fixed $\eta$, the bounded functional
$k\mapsto\langle\eta,\psi(k)\rangle$ on $L^1_+$ is integration against a
bounded function. Fubini's theorem, followed by \eqref{sc:psi}, consequently
gives
\[
\langle\eta,\psi(h*g_n)\rangle
 =\int h(y)\langle\mathsf T(y)^*\eta,\psi(g_n)\rangle\,\mathrm dy.
\]
The integrand is bounded by $C\|\eta\||h(y)|$. Weak convergence on the
right, norm convergence on the left and dominated convergence prove
\eqref{sc:point-vector}.

To locate the spectral support of $v_0$, consider the bounded function
\[
L(\xi)=\langle v_0,\mathsf T(\xi)v_0\rangle
\quad(\xi\in\overline{\mathcal W}),\qquad |L|\le C^2.
\]
It is continuous on $\overline{\mathcal W}$ and holomorphic on
$\mathcal W$.
The semigroup identity and \eqref{sc:point-vector} yield
\begin{equation}\label{sc:psi-kernel}
\langle\psi(h),\psi(k)\rangle
 =\iint\overline{h(y)}k(y')L(y^*+y')\,\mathrm dy\,\mathrm dy'
 \qquad(h,k\in L^1_+).
\end{equation}
All integrals converge absolutely by the bound on $L$.

Write $\mathfrak B(h,k)=\mu(X(h)X(k))$. If $h$ vanishes above a line
$y^0=c$ and $k$ vanishes below that line, then
\begin{equation}\label{sc:ordered-kernel}
\mathfrak B(h,k)
 =s\iint h(u)k(y)L(y-u)\,\mathrm du\,\mathrm dy.
\end{equation}
Translations reduce this to $c=0$. For that case let
$\widetilde h=\overline h\circ\theta_1\in L^1_+$. The reflection identity
gives $X(h)=s\Theta_1X(\widetilde h)$; apply
\eqref{eq:reconstructed-product} and \eqref{sc:psi-kernel}, then change variables
$y'=\theta_1u$. Since $(\theta_1u)^*=-u$, this gives
\eqref{sc:ordered-kernel}.

We next use the quarter turn to continue $L$ between the four time
directions. Let $\epsilon=(-1)^w$. The covariance of $X$ gives
\begin{equation}\label{sc:quarter-kernel}
\mathfrak B(h,k)=\epsilon\,
 \mathfrak B(h\circ R^{-1},k\circ R^{-1}).
\end{equation}
If $y_0^0>0$ and $y_0^1<0$, choose
$0<r<\frac12\min(y_0^0,-y_0^1)$ and set
\[
 h=(\pi r^2)^{-1}1_{B(0,r)},\qquad
 k=(\pi r^2)^{-1}1_{B(y_0,r)}.
\]
These functions have integral one. Both this pair and its quarter turn are
ordered across a horizontal line. Apply \eqref{sc:ordered-kernel} to the
two sides of \eqref{sc:quarter-kernel}. Their product measure has mass
one and concentrates on differences $y-u=y_0$ as $r\downarrow0$.
Continuity
of $L$ gives $L(Ry_0)=\epsilon L(y_0)$. The domain
$\mathcal W\cap R^{-1}\mathcal W$ is open, convex and contains this real
quadrant. The holomorphic identity theorem therefore gives
\begin{equation}\label{sc:quarter-L}
L(R\xi)=\epsilon L(\xi)
\qquad(\xi\in\mathcal W\cap R^{-1}\mathcal W).
\end{equation}

To use the bound on $L$, fix $\ell>0$ and put $y_\ell=(\ell,0)$.
On $R^{-m}\mathcal W$ define
$L_m(\xi)=\epsilon^mL(R^m\xi)$, $m\in\mathbb Z$.
Equation \eqref{sc:quarter-L} makes consecutive functions agree on their
overlaps, and $L_{m+4}=L_m$. If $\zeta=t+i\beta$, then
\[
R_\zeta y_\ell\in R^{-m}\mathcal W
 \quad\Longleftrightarrow\quad \cos(t+m\pi/2)>0,
\]
because the real parts of its two light-cone coordinates after applying
$R^m$ are $\ell e^{\mp\beta}\cos(t+m\pi/2)$.
These sets of angles cover $\mathbb C$. Opposite ones do not meet, and on
each other overlap the equality of consecutive $L_m$ applies. Thus
$L_m(R_\zeta y_\ell)$ defines one entire function of $\zeta$, bounded by
$C^2$. Liouville's theorem \cite[Theorem 6.16]{Dunkel} makes it constant. In particular,
\begin{equation}\label{sc:boost}
L(R_{i\beta}y_\ell)=L(y_\ell)
\qquad(\beta\in\mathbb R,\ \ell>0).
\end{equation}

Let $\varrho=\langle v_0,E_1(\cdot)v_0\rangle$.
It is a finite positive measure on $\{\omega\ge|p_2|\}$, and
\eqref{sc:boost} reads
\begin{equation}\label{sc:boost-measure}
\int e^{-\ell(\omega\cosh\beta+p_2\sinh\beta)}\,\mathrm d\varrho
 =\int e^{-\ell\omega}\,\mathrm d\varrho.
\end{equation}
On this cone
\[
\omega\cosh\beta+p_2\sinh\beta
 =\tfrac12(\omega+p_2)e^\beta+\tfrac12(\omega-p_2)e^{-\beta}\ge0.
\]
The integrand on the left of \eqref{sc:boost-measure} is at most one.
As $\beta\to+\infty$ it tends to the indicator of
$\{\omega+p_2=0\}$. With $\ell>0$ fixed, dominated convergence gives
\[
 \varrho\{\omega+p_2=0\}=\int e^{-\ell\omega}\,d\varrho.
\]
Now let $\ell\downarrow0$. Monotone convergence makes the right side
tend to the total mass of $\varrho$, so this set carries all of the
measure.
Letting $\beta\to-\infty$ gives the same assertion for
$\{\omega-p_2=0\}$. Hence $\varrho$ is carried by
$\{\omega=p_2=0\}$.

The spectral theorem now gives
$v_0\in\operatorname{Dom}H_1$ and
$\|H_1v_0\|^2=\int\omega^2\,\mathrm d\varrho=0$.
It also gives $\mathsf T(y)v_0=v_0$ for real $y^0\ge0$, since
\[
\|\mathsf T(y)v_0-v_0\|^2
 =\int|e^{-y^0\omega+iy^1p_2}-1|^2\,\mathrm d\varrho=0.
\]
Substitute this into \eqref{sc:point-vector} to conclude.
\end{proof}

\subsection{From a vacuum class to a constant random variable}

The preceding lemma is an assertion in the reconstructed Hilbert space.
The contact argument requires an assertion in $L^2(\mu)$. The $L^1$ bound
on $X$ supplies this stronger conclusion.

\begin{lemma}[Constancy of bounded local fields]\label{sc:constant-field}
Assume \textup{(S)} and $\ker H_1=\mathbb C\Omega_1$.
Every map $X$ satisfying
\eqref{sc:bounded-field}--\eqref{sc:field-actions} obeys
\[
X(h)=\mu(X(h))\,1\quad\text{in }L^2(\mu)
\qquad(h\in L^1(\mathbb R^2)).
\]
\end{lemma}

\begin{proof}
The conjugate map $X^{\mathrm c}(h)=\overline{X(\overline h)}$
satisfies the same conditions with weight $-w$. For translations and the
quarter turn this follows by commuting their operators with conjugation;
the quarter-turn phase becomes $\overline{(-i)^w}=(-i)^{-w}$.
For the reflection use that $\Theta_1$ commutes with conjugation.
The bound and positive-side measurability are preserved as well.
\Cref{sc:vacuum-classes} therefore gives
\[
[X(h)]_1,\ [X^{\mathrm c}(h)]_1\in\mathbb C\Omega_1
\qquad(h\in L^1_+).
\]
If $[\Phi]_1=a\Omega_1$ and $[\Psi]_1=b\Omega_1$, then
\eqref{sc:reconstruction} identifies $a=\mu(\Phi)$ and $b=\mu(\Psi)$.
The reflected scalar product \eqref{eq:reconstructed-product} then gives
\begin{equation}\label{sc:vacuum-factorization}
\mu((\Theta_1\Phi)\Psi)=\mu(\Theta_1\Phi)\mu(\Psi).
\end{equation}

Consider the covariance
\[
E(h,k)=\mu(\overline{X(h)}X(k))
 -\overline{\mu(X(h))}\mu(X(k)).
\]
It is sesquilinear, translation invariant, satisfies
$E(k,h)=\overline{E(h,k)}$, and obeys
\begin{equation}\label{sc:covariance-bound}
|E(h,k)|\le2C^2\|h\|_1\|k\|_1.
\end{equation}
For $h\in L^1_-$ and $k\in L^1_+$ put
$\Phi=sX^{\mathrm c}(h\circ\theta_1)$ and $\Psi=X(k)$.
Both have vacuum classes, and the reflection rule gives
$\Theta_1\Phi=\overline{X(h)}$. Thus
\eqref{sc:vacuum-factorization} says $E(h,k)=0$.
Translations and Hermitian symmetry show more generally that $E(h,k)=0$
whenever $h$ and $k$ lie on opposite sides of any line $y^0=c$.

For arbitrary $h,k\in L^1$ partition the plane into strips
$A_j=\{j\delta\le y^0<(j+1)\delta\}$.
The decompositions into these strips converge in $L^1$, and
\eqref{sc:covariance-bound} makes the resulting double series for $E(h,k)$
absolutely convergent. Every term with different strip indices vanishes by
the preceding paragraph. Hence
\[
|E(h,k)|\le2C^2\|h\|_1\sup_j\|k1_{A_j}\|_1.
\]
Set $G(a)=\int|k(a,y^1)|\,\mathrm dy^1\in L^1(\mathbb R)$.
For $M>0$,
\[
\|k1_{A_j}\|_1\le M\delta+\int G1_{\{G>M\}}.
\]
First take $M$ large and then $\delta$ small. The supremum therefore tends
to zero, so $E=0$. Its diagonal value is
$E(h,h)=\|X(h)-\mu(X(h))1\|_{L^2(\mu)}^2$, proving the lemma.
\end{proof}

\subsection{Limits of localized pairs}

Fix labels $\lambda_a=(\sigma_a,\chi_a)$, $a=1,2$, with
$\chi_a\in L^2(\mathbb R;\mathbb R)$, and put
\[
w=w(\sigma_1)+w(\sigma_2)\in\{-2,-1,0,1,2\},\qquad
s=s(\sigma_1)s(\sigma_2),\qquad
C_\lambda=2c_4^2\|\chi_1\|_2\|\chi_2\|_2.
\]
Recall that $\mathbb B_\lambda(\alpha)=\mathbb B(\alpha\otimes\lambda)$,
where $(\alpha\otimes\lambda)(x)=\alpha(y)\chi(x_3)e^\sigma$.
Define
\begin{equation}\label{sc:pair-field}
\Pi_d(z)=\mathbb B_{\lambda_1}(\varphi_{d,z})
             \mathbb B_{\lambda_2}(\varphi_{d,z}),\qquad
X_d(g)=\int g(z)\Pi_d(z)\,\mathrm dz.
\end{equation}
By \eqref{eq:pair-bound} and
$\|\varphi_{d,z}\otimes\lambda_a\|\le\sqrt2\|\chi_a\|_2$,
the map $\Pi_d$ is continuous with
\begin{equation}\label{sc:pair-field-bound}
\|\Pi_d(z)\|_{L^2(\mu)}\le C_\lambda,\qquad
\|X_d(g)\|_{L^2(\mu)}\le C_\lambda\|g\|_1.
\end{equation}
In particular, the integral is well defined for every $g\in L^1$.

\begin{lemma}[Localized pair fields]\label{sc:pair-limits}
Assume $(G)$, $(T)$ and $(Q_{12})$. The translation and quarter-turn
identities below and the mean formula hold under these assumptions.
If $(R_1)$ holds as well, so does the reflection identity:
\begin{align}
&U(v)X_d(g)=X_d(\tau_vg),\qquad
 \Theta_1X_d(g)=sX_d(\overline g\circ\theta_1),\qquad
 \hat RX_d(g)=(-i)^wX_d(g\circ R^{-1}),
 \label{sc:pair-covariance}\\
&\mu(X_d(g))=\kappa_d\int g,\qquad
 \kappa_d=\mu(\Pi_d(0)),\qquad |\kappa_d|\le C_\lambda.
 \label{sc:pair-mean}
\end{align}
If $w\ne0$, then $\kappa_d=0$.
For the remaining assertions assume also $(R_1)$ and $(RP_1)$.
If $g$ vanishes on $\{z^0<d\}$, then
$X_d(g)\in L^2(\Sigma_1^+,\mu)$.

For the fixed labels $\lambda_1,\lambda_2$, every sequence
$d_n\downarrow0$ has a subsequence, common to all $g\in L^1$, along
which $X_{d_n}(g)$ converges weakly in $L^2(\mu)$.
Its limit satisfies
\eqref{sc:bounded-field}--\eqref{sc:field-actions} with constant
$C_\lambda$, sign $s$ and weight $w$.

If in addition \textup{(S)} holds and
$\ker H_1=\mathbb C\Omega_1$, then, for all
$g\in L^1$ and $\Phi\in L^2(\mu)$,
\begin{equation}\label{sc:pair-factorization}
\mu(\Phi X_d(g))-\mu(\Phi)\mu(X_d(g))\longrightarrow0
\qquad(d\downarrow0).\qedhere
\end{equation}
\end{lemma}

\begin{proof}
The field actions and \eqref{sc:disc-actions} give
\[
U(v)\Pi_d(z)=\Pi_d(z+v),\qquad
\Theta_1\Pi_d(z)=s\Pi_d(\theta_1z),\qquad
\hat R\Pi_d(z)=(-i)^w\Pi_d(Rz).
\]
The reflection formula uses that each $\varphi_{d,z}$ is real; the
quarter-turn formula uses
$R\cdot(\alpha\otimes\lambda)=(-i)^{w(\sigma)}
(\alpha\circ R^{-1})\otimes\lambda$.
Pass these operators through the integral in \eqref{sc:pair-field},
conjugating the scalar coefficient for $\Theta_1$, and change variables.
This proves \eqref{sc:pair-covariance}.
Translation invariance makes $\mu(\Pi_d(z))$ independent of $z$, proving
\eqref{sc:pair-mean}; the bound follows from
\eqref{sc:pair-field-bound}. At $z=0$ the quarter-turn identity gives
$\kappa_d=(-i)^w\kappa_d$. Since $w$ is a pair weight,
$(-i)^w\ne1$ for $w\ne0$, so $\kappa_d=0$ in that case.
For $z^0\ge d$ the disc $B(z,d)$ lies in the positive half-plane.
Each field in $\Pi_d(z)$ is then measurable on the positive side and
belongs to every finite $L^p(\mu)$; their product belongs to
$L^2(\Sigma_1^+,\mu)$. The integral stays in this closed subspace.

For the weak limits, choose a countable dense set $\{g_j\}$ in $L^1$,
for example finite sums of dyadic-square indicators with coefficients
in $\mathbb Q+i\mathbb Q$.
The bound \eqref{sc:pair-field-bound} and successive weak
subsequence extractions give a single subsequence along which every
$X_{d_n}(g_j)$ converges weakly. For arbitrary $g\in L^1$ and
$\Psi\in L^2(\mu)$,
\[
\begin{split}
|\langle\Psi,X_{d_n}(g)-X_{d_m}(g)\rangle|
\le{}&2C_\lambda\|\Psi\|\|g-g_j\|_1\\
&+|\langle\Psi,X_{d_n}(g_j)-X_{d_m}(g_j)\rangle|.
\end{split}
\]
The scalar pairings are therefore Cauchy. Their limits define, by the
Riesz theorem, a vector $X(g)$ with
$\|X(g)\|\le C_\lambda\|g\|_1$. Linearity and the covariance identities
pass to the limit. This is also true for the antiunitary $\Theta_1$,
since $\langle\Psi,\Theta_1v\rangle
=\overline{\langle\Theta_1\Psi,v\rangle}$.

It remains to check positive-side measurability of $X(g)$ when
$g\in L^1_+$. If $g$ vanishes on $\{z^0<\delta\}$ for some $\delta>0$,
then $X_{d_n}(g)\in L^2(\Sigma_1^+,\mu)$ once $d_n<\delta$.
A closed subspace of a Hilbert space is weakly closed, so the same holds
for $X(g)$. For general $g\in L^1_+$, approximate it by
$g1_{\{z^0\ge\delta\}}$ in $L^1$ and use the bound on $X$.
Thus all the conditions on the limiting field hold.

Finally, suppose \eqref{sc:pair-factorization} fails for fixed $g,\Phi$.
There are $\delta>0$ and $d_n\downarrow0$ for which the modulus of its
left side is at least $\delta$. Pass to the subsequence just constructed.
\Cref{sc:constant-field} makes its limiting field constant:
$X(g)=\mu(X(g))1$. Weak convergence, paired with $\overline\Phi$ and
with $1$, then sends that same left side to zero, a contradiction.
\end{proof}

\subsection{Pairs at contact}

For arbitrary $f_1,f_2\in\mathfrak h$ define
\begin{equation}\label{sc:Y-definition}
Y_d(f_1,f_2)=\int
 \mathbb B(\varphi_{d,z}f_1)\mathbb B(\varphi_{d,z}f_2)\,\mathrm dz.
\end{equation}
The next lemma compares this contact pair with the fields just studied.

\begin{lemma}[Contact-pair estimates]\label{sc:contact-pairs}
Assume \textup{(G)}. The integral \eqref{sc:Y-definition} exists in $L^2(\mu)$,
is bilinear in its entries, and obeys
\begin{equation}\label{sc:Y-bound}
\|Y_d(f_1,f_2)\|_{L^2(\mu)}\le c_4^2\|f_1\|\|f_2\|.
\end{equation}
If $f_a=\alpha_a\otimes\lambda_a$ with
$\alpha_a\in C_{\mathrm c}(\mathbb R^2)$, then
\begin{equation}\label{sc:Y-X}
\|Y_d(f_1,f_2)-X_d(\alpha_1\alpha_2)\|_{L^2(\mu)}\longrightarrow0.
\end{equation}
Under $(T)$ and $(Q_{12})$, for arbitrary $f_1,f_2\in\mathfrak h$
and $t\in\mathbb R$,
\begin{equation}\label{sc:Y-mean-rotation}
\mu(Y_d(R_t\cdot f_1,R_t\cdot f_2))-\mu(Y_d(f_1,f_2))
 \longrightarrow0.
\end{equation}
Under \textup{(S)} and $\ker H_1=\mathbb C\Omega_1$, for arbitrary
$f_1,f_2\in\mathfrak h$ and $\Phi\in L^2(\mu)$,
\begin{equation}\label{sc:Y-factorization}
\mu(\Phi Y_d(f_1,f_2))-\mu(\Phi)\mu(Y_d(f_1,f_2))
 \longrightarrow0.
\end{equation}
All limits are as $d\downarrow0$.
\end{lemma}

\begin{proof}
The integrand in \eqref{sc:Y-definition} is norm continuous. By
\eqref{eq:pair-bound}, Cauchy--Schwarz in $z$, and
\eqref{sc:discs},
\[
\int\|\mathbb B(\varphi_{d,z}f_1)
             \mathbb B(\varphi_{d,z}f_2)\|_{L^2(\mu)}\,\mathrm dz
 \le c_4^2\int\|\varphi_{d,z}f_1\|\|\varphi_{d,z}f_2\|\,\mathrm dz
 \le c_4^2\|f_1\|\|f_2\|.
\]
This proves existence, bilinearity and \eqref{sc:Y-bound}. It also permits
pairing the integral with any $\Phi\in L^2(\mu)$ as in
\eqref{sc:integral-pairing}.

For \eqref{sc:Y-X} let
$\omega_a(d)=\sup_{|y-z|\le d}|\alpha_a(y)-\alpha_a(z)|$.
Uniform continuity gives $\omega_a(d)\to0$. Moreover,
\[
\|\varphi_{d,z}\alpha_a-\alpha_a(z)\varphi_{d,z}\|_2\le\omega_a(d),
\qquad
\|\varphi_{d,z}\alpha_a\|_2,
\ \|\alpha_a(z)\varphi_{d,z}\|_2\le\|\alpha_a\|_\infty.
\]
Expand the difference of the two products of fields, changing one factor
at a time, and apply \eqref{eq:pair-bound}. Its norm is at most
\[
C_\lambda\bigl(\omega_1(d)\|\alpha_2\|_\infty
             +\|\alpha_1\|_\infty\omega_2(d)\bigr).
\]
For $d\le1$ it vanishes outside the fixed compact one-neighborhood $K$
of $\operatorname{supp}\alpha_1\cup\operatorname{supp}\alpha_2$.
Integration therefore bounds the left side of \eqref{sc:Y-X} by $|K|$
times this expression, proving the limit.

The bilinear forms on the left of \eqref{sc:Y-mean-rotation} are bounded
uniformly in $d$ by \eqref{sc:Y-bound}; those on the left of
\eqref{sc:Y-factorization} have the same property for fixed $\Phi$.
It suffices to prove each limit on finite sums of plates with continuous
compactly supported plane profiles. These form a dense subspace of
$\mathfrak h$: first use totality of plates, then approximate each plane
profile in $L^2$ by a continuous function of compact support.
The extension follows, for any uniformly bounded family $T_d$, from
\[
|T_d(f_1,f_2)|\le |T_d(f'_1,f'_2)|
 +M\bigl(\|f_1-f'_1\|\|f_2\|+\|f'_1\|\|f_2-f'_2\|\bigr).
\]

We may therefore take $f_a=\alpha_a\otimes\lambda_a$ with continuous
compactly supported $\alpha_a$. The plate rotation formula,
\eqref{sc:Y-X}, and \eqref{sc:pair-mean} give
\[
\begin{split}
&\mu(Y_d(R_t\cdot f_1,R_t\cdot f_2))-\mu(Y_d(f_1,f_2))\\
&\qquad=(e^{-iwt}-1)\kappa_d\int\alpha_1\alpha_2+o(1).
\end{split}
\]
If $w=0$, the phase difference vanishes. If $w\ne0$, then
$\kappa_d=0$ by \cref{sc:pair-limits}. No convergence of
$\kappa_d$ as $d\downarrow0$ is needed. This proves
\eqref{sc:Y-mean-rotation}. Finally, \eqref{sc:Y-X} changes the expression
in \eqref{sc:Y-factorization} by $o(1)$ when $Y_d$ is replaced by
$X_d(\alpha_1\alpha_2)$. The latter expression tends to zero by
\eqref{sc:pair-factorization}. Density completes both proofs.
\end{proof}

\subsection{Removing the disjointness conditions}

We now state the precise extension used in the proof of the main theorem.
For $N\ge0$ and $t\in\mathbb R$ let
\[
D_N^t(f_1,\ldots,f_N)
 =S_N(R_t\cdot f_1,\ldots,R_t\cdot f_N)-S_N(f_1,\ldots,f_N),
\qquad S_0=1.
\]
The Gaussian bounds imply that each $D_N^t$ is bounded, symmetric and
$N$-linear, and $D_0^t=D_1^t=0$.

\begin{proposition}[From disjoint supports to all entries]
\label{sc:all-entries}
Assume \textup{(S)} and $\ker H_1=\mathbb C\Omega_1$.
Suppose that, for every $N\ge2$ and $t\in\mathbb R$,
\[
D_N^t(f_1,\ldots,f_N)=0
\]
whenever $f_a\in\mathfrak h_{K_a}$ for compact pairwise disjoint
$K_1,\ldots,K_N\subset\mathbb R^2$.
Then $D_N^t=0$ for every $N\ge0$ and every $t\in\mathbb R$.
Consequently $(R_t)_*\mu=\mu$ for all $t$.
\end{proposition}

\begin{proof}
Fix $t$, write $D_N=D_N^t$, and suppose inductively that $D_{N-2}=0$.
The cases $N=0,1$ have already been noted. For a collection $E$ of
two-element subsets of $\{1,\ldots,N\}$, let $\mathcal P(E)$ mean that
$D_N$ vanishes on entries $f_a\in\mathfrak h_{K_a}$ whenever the $K_a$
are compact and $K_a\cap K_b=\varnothing$ for every $\{a,b\}\in E$.
The assumed separated-support invariance gives $\mathcal P(E)$ when
$E$ contains every pair. We show that any one pair can be removed.

Suppose $\mathcal P(E)$ holds and $e=\{a,b\}\in E$.
Fix compact sets $K_c$ satisfying all the disjointness conditions in
$E\setminus\{e\}$ and entries $f_c\in\mathfrak h_{K_c}$.
With the other entries fixed, regard $D_N$ as a bounded bilinear form
$\beta$ on $\mathfrak h_{K_a}\times\mathfrak h_{K_b}$.
If its two arguments are further restricted over disjoint compact sets
$A,A'$, then the sets $K_a\cap A$, $K_b\cap A'$ and the remaining $K_c$
satisfy every condition in $E$. Hence $\mathcal P(E)$ makes this value of
$\beta$ zero. \Cref{sc:localization} applies and gives, for all $d>0$,
\begin{equation}\label{sc:pair-removal-start}
D_N(f_1,\ldots,f_N)
 =\int\beta(\varphi_{d,z}f_a,\varphi_{d,z}f_b)\,\mathrm dz.
\end{equation}

Put
\[
\Phi=\prod_{c\ne a,b}\mathbb B(f_c),\qquad
\Phi'=\prod_{c\ne a,b}\mathbb B(R_t\cdot f_c),
\]
with empty products equal to one. Both belong to $L^2(\mu)$ by Gaussian
domination. Using \eqref{sc:disc-actions} in the rotated term of
\eqref{sc:pair-removal-start}, and then using $R_tz$ as the integration
variable in that term, yields
\begin{equation}\label{sc:pair-removal-Y}
D_N(f_1,\ldots,f_N)
 =\mu(Y_d(R_t\cdot f_a,R_t\cdot f_b)\Phi')
  -\mu(Y_d(f_a,f_b)\Phi).
\end{equation}
Each term is integrable before this change of variables, by the integral
estimate in \cref{sc:contact-pairs} and Cauchy--Schwarz in $L^2(\mu)$.

The induction hypothesis gives $\mu(\Phi')=\mu(\Phi)$. Subtract and add
the corresponding products of means in \eqref{sc:pair-removal-Y}:
\[
\begin{split}
D_N(f_1,\ldots,f_N)={}&
 \bigl[\mu(Y_d(R_t\cdot f_a,R_t\cdot f_b)\Phi')
  -\mu(Y_d(R_t\cdot f_a,R_t\cdot f_b))\mu(\Phi')\bigr]\\
&-\bigl[\mu(Y_d(f_a,f_b)\Phi)
  -\mu(Y_d(f_a,f_b))\mu(\Phi)\bigr]\\
&+\mu(\Phi)\bigl[\mu(Y_d(R_t\cdot f_a,R_t\cdot f_b))
  -\mu(Y_d(f_a,f_b))\bigr].
\end{split}
\]
As $d\downarrow0$, the first two brackets vanish by
\eqref{sc:Y-factorization}; the third vanishes by
\eqref{sc:Y-mean-rotation}. The left side is independent of $d$, so it is
zero. This proves $\mathcal P(E\setminus\{e\})$.

Remove the pairs successively to obtain $\mathcal P(\varnothing)$.
Thus $D_N$ vanishes whenever its entries each lie over a compact subset
of the plane, without any condition on intersections. These entries are
dense in $\mathfrak h$, since $1_{\{|y|\le n\}}f\to f$ in $L^2$.
Continuity gives $D_N=0$ on $\mathfrak h^N$.
Induction from $D_0=0$ for even $N$ and from $D_1=0$ for odd $N$ proves
the assertion for all $N$ and $t$.

In particular all moments of real Schwartz fields are invariant under
$R_t$. The moment-determination consequence of Gaussian domination then
gives $(R_t)_*\mu=\mu$.
\end{proof}



\section{Ergodic components and the full Euclidean group}\label{sec:components}

We complete the proof of \cref{thm:main}. By
\cref{lem:transport}, the law has (H) and belongs to $\mathcal C_9$.
The ergodic-component arguments in
\cite[Sections~7.3--7.5]{QEDCurrent} apply under these hypotheses.
For the translation-ergodic decomposition
\[
 \mu=\int\mu_\xi\,d\nu(\xi),
\]
they give, outside a single $\nu$-null set,
\begin{equation}\label{eq:components}
 \mu_\xi\in\mathcal C_9,\qquad
 \mu_\xi\text{ satisfies }(G),\qquad
 \mu_\xi\text{ is invariant under }\R^3\rtimes\Ocub,\qquad
 \ker H_1(\mu_\xi)=\C\Omega_1.
\end{equation}
In particular each such component has (H). Neither this component
theorem nor its hypotheses require closedness.

Fix one component with these properties. The cone theorem gives (S),
and \cref{prop:separated} gives invariance under $R_t$ on
tests over pairwise disjoint compact planar sets, at every order.
\Cref{sc:all-entries} therefore gives
$(R_t)_*\mu_\xi=\mu_\xi$ for every $t\in\R$.

The orthogonal symmetries of $\mu_\xi$ form a group containing
$\Ocub$ and all $R^{12}_t$. Conjugation by a cyclic permutation of
the axes gives all $R^{23}_t$. Every matrix in $SO(3)$ is a product
$R^{12}_{t_1}R^{23}_{t_2}R^{12}_{t_3}$: choose the first two rotations
to send $e_3$ to its image under the matrix; the remaining factor
fixes $e_3$ and is a rotation in its orthogonal plane. Since the cubic
group also contains a reflection, this proves invariance under all
of $O(3)$. It holds for every $g$ on the same full-measure set of
components; no intersection over an uncountable family of null sets
is being taken.

For any measurable set $A$ of distributions, the action is measurable
because its evaluation on a test $f$ is evaluation on $g^{-1}\cdot f$.
Consequently
\[
 (g_*\mu)(A)=\int\mu_\xi(g^{-1}A)\,d\nu(\xi)
            =\int\mu_\xi(A)\,d\nu(\xi)=\mu(A).
\]
Translation invariance gives the corresponding statement for every
Euclidean isometry. Finally \eqref{eq:moments} and the action on
$\B$ give moment covariance on all $L^2$ tests. This proves
\cref{thm:main}.

\section{Diagonal positivity, the cone and closedness}\label{sec:examples}

The diagonal hypothesis cannot be replaced by the cone alone. We give
explicit negative reflection forms for two laws that satisfy (H) and
all ordered cone conditions $(C_{ji})$ but are not closed. We then
give an invariant law positive at every plane for which closedness
also fails. Together these examples distinguish the hypothesis of
\cref{thm:main} from both the cone and the Bianchi identity.

For real Schwartz vector tests and a bounded measurable, even, real
symmetric positive-semidefinite matrix multiplier $M$, write
\[
 C_M(b,b')=(2\pi)^{-3}\int\widehat b(p)^\dagger M(p)\widehat b'(p)\,dp,
 \qquad \widehat b(p)=\int e^{-ip\cdot x}b(x)\,dx.
\]
Let $\gamma_M$ be the centered Gaussian law with characteristic
functional $\gamma_M(e^{i\B(b)})=e^{-C_M(b,b)/2}$.
The Gaussian covariance criterion in
\cite[Appendix~D, ``Gaussian laws and reflection positivity'']{QEDConsequences}
gives existence, and $M\le I$ gives (G) by Plancherel's identity.
The phases of a common translation cancel in $C_M$, so all these
laws are translation invariant. For $g\in O(3)$, changing variables in
the Fourier integral shows that $M(gp)=gM(p)g^T$ almost everywhere implies
invariance under $g$. Closedness is equivalent to
$M(p)p=0$ almost everywhere: indeed,
\[
 C_M(\nabla u,\nabla u)
 =(2\pi)^{-3}\int|\widehat u(p)|^2p^TM(p)p\,dp.
\]
If $M(p)p=0$ almost everywhere, the displayed integral vanishes for
every $u$. Conversely, closedness makes it vanish for a Gaussian
$u$, whose Fourier transform never vanishes; positivity of $M$ then
gives $M(p)p=0$ almost everywhere.
We also use the Gaussian reflection criterion from the same reference:
if $C_M$ is reflection invariant, its law is reflection positive
precisely when $C_M(r\cdot b,b)\ge0$ for every real smooth test
compactly supported in the positive half-space.

\subsection{A diagonal Maxwell covariance}

For $m\ge0$, define
\[
 M^\oplus_{ab}(p)=\delta_{ab}\frac{|p|^2-p_a^2}{|p|^2+m^2}
 \quad(p\ne0),\qquad M^\oplus(0)=0,
 \qquad \gamma_m^\oplus=\gamma_{M^\oplus}.
\]
This multiplier lies between zero and $I$ and obeys
$M^\oplus(gp)=gM^\oplus(p)g^T$ for signed permutations $g$.
Thus the law has (G), translation invariance and cubic symmetry.
To verify coordinate positivity, write $p=(p_1,k)$ and
$\omega(k)=(|k|^2+m^2)^{1/2}$. At the single point $k=0$ when
$m=0$, set $\omega(0)=1$; this does not change any integral.
The partial fraction decomposition is
\[
 M^\oplus(p)=\operatorname{diag}(0,1,1)
  +\frac{\operatorname{diag}(|k|^2,-p_2^2-m^2,-p_3^2-m^2)}
             {p_1^2+\omega(k)^2}.
\]
The residue matrix in \cref{lem:example-time-kernel} is
$\operatorname{diag}(|k|^2,p_2^2+m^2,p_3^2+m^2)\ge0$.
That lemma gives $(RP_1)$, and the spectral measures of all
one-field vectors are supported on
$\{(\omega,k):\omega=(|k|^2+m^2)^{1/2}\}$.
The round cone $\{\omega\ge|k|\}$ is closed, contains zero and
is stable under addition. Lemma~\ref{lem:example-gaussian-cone}
therefore puts the whole reconstructed spectrum in this cone.
Cubic symmetry gives all coordinate reflection positivities and
all ordered cones $(C_{ji})$. Thus the law has (H).
It is not closed: $p^TM^\oplus(p)p>0$ whenever $p$ has at least
two nonzero components. The following test identifies its failure
at every diagonal mirror.

\begin{proposition}\label{prop:diagonal-maxwell}
For every $m\ge0$, the law $\gamma_m^\oplus$ fails reflection
positivity at every plane $x_i=\sigma x_l+c$, $i<l$, $\sigma=\pm1$,
from either side. For a fixed plane and side one can use the same
test forms for all $m\ge0$.
\end{proposition}

\begin{proof}
First take $c=0$. Fix signs $\sigma,\nu\in\{1,-1\}$ and put
\[
 n=\frac{\nu}{\sqrt2}(e_i-\sigma e_l),\qquad
 H_+=\{n\cdot x>0\},\qquad r=r_{n,0}.
\]
The reflection exchanges $e_i$ and $\sigma e_l$ and fixes the third
basis vector. Its action on tests is $r\cdot b=-r\,b\circ r$.
If $C_a(u',u)=C_{M^\oplus}(u'e_a,ue_a)$, the multiplier identities
$M^\oplus_{ll}(rp)=M^\oplus_{ii}(p)$ and
$M^\oplus_{ii}(rp)=M^\oplus_{ll}(p)$ give, for $b=ue_i+ve_l$,
\begin{equation}\label{eq:reflected-diagonal}
 C_{M^\oplus}(r\cdot b,r\cdot b)=C_{M^\oplus}(b,b),\qquad
 C_{M^\oplus}(r\cdot b,b)=-2\sigma C_i(v\circ r,u).
\end{equation}

Between disjoint supports the scalar covariance has kernel
\[
 K_i(x)=(\partial_i^2-m^2)\frac{e^{-m|x|}}{4\pi|x|}.
\]
Indeed, write its multiplier as
$1-(p_i^2+m^2)/(|p|^2+m^2)$; the identity term vanishes between
disjoint supports, and Fourier inversion of the Yukawa potential
followed by integration by parts gives this formula. It remains
valid for $m=0$ by dominated convergence, since $|p|^{-2}$ is locally
integrable in three dimensions. With $\rho=|x|$ and $s=m\rho$,
direct differentiation yields
\begin{equation}\label{eq:negative-kernel}
 K_i(x)=\frac{e^{-s}}{4\pi\rho^3}
 \left[(s^2+3s+3)\frac{x_i^2}{\rho^2}-(s^2+s+1)\right].
\end{equation}
In particular, since $s^2+3s+3\le3(s^2+s+1)$,
\begin{equation}\label{eq:negative-region}
 K_i(x)\le-\frac{e^{-m\rho}}{16\pi\rho^3}<0
 \quad\text{if }x_i^2\le\rho^2/4.
\end{equation}

Choose $\lambda>0$ and a nonzero $\varphi\in C_c^\infty(\R^3)$
with $\varphi\ge0$ and support in $\{|x|\le\lambda/3\}$.
Set
\[
 y_1=\nu\sigma\lambda e_l,\qquad
 u(x)=\varphi(x+y_1),\qquad
 v(x)=-\sigma\varphi(rx-y_1),\qquad b=ue_i+ve_l.
\]
The two supports lie in balls centered at $-\nu\sigma\lambda e_l$
and $\nu\lambda e_i$. Both centers have $n$-coordinate
$\lambda/\sqrt2$, so both closed balls lie in $H_+$.
The form $f=*^{-1}b$ is therefore an admissible test in
\eqref{eq:rp}; it does not depend on $m$.

By \eqref{eq:reflected-diagonal}, the common variance is
$V=2C_i(\varphi,\varphi)$ and the reflected covariance is
\[
 J=2\iint\varphi(y)\varphi(y')
             K_i(y-y'+2\lambda e_l)\,dy\,dy'.
\]
To obtain the displayed sign of the displacement from $2y_1$, use
evenness of $K_i$ and interchange $y,y'$ if necessary.
On the integration region,
\[
 |(y-y'+2\lambda e_l)_i|\le2\lambda/3,
 \qquad4\lambda/3\le|y-y'+2\lambda e_l|\le8\lambda/3.
\]
Thus \eqref{eq:negative-region} applies throughout, and
\begin{equation}\label{eq:J-negative}
 J\le-c\lambda^{-3}e^{-8m\lambda/3}\left(\int\varphi\right)^2<0
\end{equation}
with an absolute constant $c>0$.

For Gaussian data $f,-f$ and coefficients $1,-1$, expansion of the
four terms in \eqref{eq:rp} gives
\begin{equation}\label{eq:gaussian-negative}
 \mathcal Q_{\gamma_m^\oplus}(n,0;(1,-1);(f,-f))
       =4e^{-V}\sinh J<0.
\end{equation}
One may instead use tests $0,f$ and coefficients $1,-e^{V/2}$;
their form is $e^J-1<0$. This proves failure at the chosen plane
and side. Translation of the tests carries these same values to
every parallel plane, because
$r_{n,c}=\tau_{cn}r_{n,0}\tau_{-cn}$ and the law is translation
invariant. The choice of $\nu$ includes both sides.
\end{proof}

For radial $\varphi$, the witness also has an exact expression.
Define
\[
 c_\varphi(m)=\int\varphi(y)\frac{\sinh(m|y|)}{m|y|}\,dy,
\]
with quotient one at $m=0$. Then
\begin{align}
 J&=-c_\varphi(m)^2
       \frac{(4m^2\lambda^2+2m\lambda+1)e^{-2m\lambda}}
            {16\pi\lambda^3},\label{eq:radial-J}\\
 V&=\frac43(2\pi)^{-3}\int|\widehat\varphi(p)|^2
                  \frac{|p|^2}{|p|^2+m^2}\,dp.
 \label{eq:radial-V}
\end{align}
To verify the first formula, a smooth solution of $\Delta G=m^2G$
near a ball has spherical mean
$G(0)\sinh(m\rho)/(m\rho)$. Green's formula gives the radial equation
$S''+2\rho^{-1}S'=m^2S$, with $S(0)=G(0)$; solving for $\rho S$
gives this mean. Radial integration therefore gives
$\int\varphi G=c_\varphi(m)G(0)$. Apply it twice to the two
variables in $J$, using $\Delta K_i=m^2K_i$ away from the origin.
Then $J=2c_\varphi(m)^2K_i(2\lambda e_l)$, which is
\eqref{eq:radial-J}. Symmetry among the three Fourier coordinates
gives \eqref{eq:radial-V}. At $m=0$ these reduce to
\[
 J=-\frac{(\int\varphi)^2}{16\pi\lambda^3},\qquad
 V=\frac43\|\varphi\|_2^2.
\]

\subsection{A mixture of white noises}

Let $D_1,D_2,D_3\in(0,1]$ be pairwise distinct. For each permutation
$\pi\in\mathfrak S_3$, let $\gamma_{\pi D}$ be the white-noise law
with variance $D_{\pi^{-1}(a)}$ in component $a$ of $\B$, and put
\[
 m_1=\frac16\sum_{\pi\in\mathfrak S_3}\gamma_{\pi D}.
\]
Each summand satisfies (G), is translation invariant and is
invariant under coordinate reflections. The cubic group permutes
the six summands. Opposite open half-spaces are independent for
each white noise, since their covariance vanishes. Consequently,
for positive-side $\Phi,\Psi\in L^2(m_1)$,
\[
 m_1((\Theta_j\Phi)\Psi)
 =\frac16\sum_{\pi\in\mathfrak S_3}
       \overline{\gamma_{\pi D}(\Phi)}\gamma_{\pi D}(\Psi).
\]
The inclusion $L^2(m_1)\subset L^2(\gamma_{\pi D})$ follows from
$m_1\ge\gamma_{\pi D}/6$. The displayed formula proves coordinate
reflection positivity. It also identifies a reconstructed class
by its six expectations, each unchanged by translations.
Hence every reconstructed translation is the identity, so
$H_j=0$ and all transverse momenta vanish. This proves (H) and
all ordered cones. Moreover
\[
 m_1(\B(\nabla u)^2)
   =\frac{D_1+D_2+D_3}{3}\|\nabla u\|_2^2>0
\]
for every nonzero real Schwartz function $u$, so $m_1$ is not closed.
Although the individual white noises have no correlations between
opposite open half-spaces, their mixture need not be reflection
positive at a reflection that permutes the variances.

Fix the plane, side and reflection from
\cref{prop:diagonal-maxwell}. Take a nonzero real
$\psi\in C_c^\infty(H_+)$ and set $*f_1=\psi e_i$,
$*f_2=\psi e_l$. Let $\mathfrak t$ exchange $i,l$ and write
\[
 \beta_a=e^{-D_a\|\psi\|_2^2/2},\qquad
 Z_\pi=\beta_{\pi^{-1}(i)}-\beta_{\pi^{-1}(l)}.
\]
For tests on the positive side, the white-noise characteristic
functional factors between the two half spaces. Reflection permutes
the weights by $\mathfrak t$, so its form for coefficients $(1,-1)$
is
\begin{align}
 \mathcal Q_{m_1}(n,0;(1,-1);(f_1,f_2))
 &=\frac16\sum_\pi\overline{Z_{\mathfrak t\pi}}Z_\pi\notag\\
 &=-\frac16\sum_\pi Z_\pi^2
  =-\frac13\sum_{a<c}(\beta_a-\beta_c)^2<0.
 \label{eq:mixture-negative}
\end{align}
Here $Z_{\mathfrak t\pi}=-Z_\pi$, and the preimages of $i,l$ run
through all six ordered pairs of distinct indices. Distinct variances
make the last inequality strict. The calculation works for both
values of $\nu$; translation gives the same failure at every parallel
plane. Thus neither $m_1$ nor $\gamma_m^\oplus$ belongs to
$\mathcal C_9$.

\subsection{An invariant nonclosed law}

The preceding examples fail the extra positivity hypothesis of
\cref{thm:main}. The following example satisfies it, and shows
that the theorem applies to nonclosed laws.

\begin{proposition}[An invariant nonclosed law]\label{prop:nonclosed}
For $m>0$, let $M(p)=I-pp^T/(|p|^2+m^2)$. The Gaussian law
$\gamma_M$ has (H), is invariant under every Euclidean isometry,
and is reflection positive at every affine plane. It is not closed.
\end{proposition}

\begin{proof}
On $p^\perp$ the multiplier is $I$, and on $\R p$ its eigenvalue is
$m^2/(|p|^2+m^2)$. Thus $0\le M\le I$; it is measurable, even and
real symmetric. The multiplier criteria give the Gaussian law,
translation invariance and (G). Since
\[
 M(gp)=gM(p)g^T\quad(g\in O(3)),\qquad
 M(p)p=\frac{m^2}{|p|^2+m^2}p,
\]
the law is orthogonally invariant and not closed.

For coordinate reflection positivity, split $p=(p_1,p_\perp)$ and
put $\omega=(|p_\perp|^2+m^2)^{1/2}$.
Lemma~\ref{lem:example-time-kernel} applies to
\[
 M(p)=I+\frac{N(p)}{p_1^2+\omega^2},\qquad N(p)=-pp^T.
\]
Its polynomial part is constant, its numerator is polynomial in
$p_1$, and $\omega\ge m>0$. It obeys the required reflection
symmetry $\theta_1M(-p_1,p_\perp)\theta_1=M(p)$.
The residue matrix in that criterion, with
$\mathsf p=(-i\omega,p_\perp)$, is
\[
 A=-\theta_1N(-i\omega,p_\perp)
   =\overline{\mathsf p}\,\mathsf p^T,\qquad
 v^\dagger Av=|\mathsf p^Tv|^2\ge0.
\]
The criterion therefore gives $(RP_1)$. Orthogonal and translation
invariance transport it to every affine plane, including both sides.
In particular all three coordinate conditions hold, so the law has
(H).
\end{proof}

Closedness and diagonal positivity consequently play different roles.
The cone and (H) alone do not give diagonal positivity, as
\eqref{eq:gaussian-negative} and \eqref{eq:mixture-negative} show.
Diagonal positivity does give rotation invariance under (H), by
\cref{thm:main}, without forcing closedness, by
\cref{prop:nonclosed}. The elimination of the extra two-point
frequency by the Bianchi identity, recalled in
\cref{eq:bianchi} and proved in \cite[Section~8.5]{QEDCurrent},
still uses $dF=0$. The examples above show why a conclusion of
diagonal positivity from (H) and the cone requires an additional
hypothesis.

\appendix
\section{Gaussian reflection forms and their spectral support}
\label{app:gaussian-examples}

We record the two covariance calculations used in
\cref{sec:examples}. Write $p=(p_1,k)$, with $k\in\R^2$, and regard
$\theta_1$ as the matrix $\operatorname{diag}(-1,1,1)$.
In the coordinate reconstruction, write
$V_1(y)=\int e^{ik\cdot y}\,dE_1(\omega,k)$ for the transverse
translation, with $y\in e_1^\perp$.

\begin{lemma}[A time-kernel formula]\label{lem:example-time-kernel}
Let $M$ be bounded, measurable, even, real symmetric and positive
semidefinite, and suppose
$\theta_1M(-p_1,k)\theta_1=M(p_1,k)$. Suppose that, for almost every
$k$, there are matrix polynomials $Q_k,N_k$ and a measurable
$\omega(k)>0$ such that
\[
 M(p_1,k)=Q_k(p_1)+\frac{N_k(p_1)}{p_1^2+\omega(k)^2}.
\]
The coefficients of these polynomials are assumed measurable in $k$.
Put $A(k)=-\theta_1N_k(-i\omega(k))$. For a real
$b\in C_c^\infty(\{x_1>0\};\R^3)$ define
\[
 L_b(k)=\int_0^\infty e^{-\omega(k)s}\widehat b_\perp(s,k)\,ds,
 \qquad
 \widehat b_\perp(s,k)=\int_{\R^2}e^{-ik\cdot x_\perp}b(s,x_\perp)\,dx_\perp.
\]
Then, for $t\ge0$ and $y\in e_1^\perp$,
\begin{equation}\label{eq:example-time-kernel}
 C_M(\theta_1\cdot b,\tau_{te_1+y}\cdot b)
 =(2\pi)^{-2}\int_{\R^2}e^{-ik\cdot y}
       \frac{e^{-\omega(k)t}}{2\omega(k)}
       L_b(k)^\dagger A(k)L_b(k)\,dk.
\end{equation}
If $A(k)$ is Hermitian and positive semidefinite almost everywhere,
$\gamma_M$ is reflection positive at $x_1=0$. Under this positivity
hypothesis, the spectral measure of $[\B(b)]_1$ is the image under
$k\mapsto(\omega(k),-k)$ of
\[
 (2\pi)^{-2}\frac{L_b(k)^\dagger A(k)L_b(k)}{2\omega(k)}\,dk.
\]
\end{lemma}

\begin{proof}
For $h\in C_c^\infty((0,\infty))$ and a polynomial $P$,
Fourier inversion and the inverse transform of
$(p_1^2+\omega^2)^{-1}$ give
\begin{align*}
 \frac1{2\pi}\int\widehat h(p_1)P(p_1)\,dp_1&=0,\\
 \frac1{2\pi}\int\frac{\widehat h(p_1)P(p_1)}{p_1^2+\omega^2}\,dp_1
 &=\frac{P(-i\omega)}{2\omega}
                   \int_0^\infty e^{-\omega u}h(u)\,du.
\end{align*}
For the first identity, derivatives of $h$ vanish at zero. For the
second, the case $P=1$ follows from the inverse transform
$e^{-\omega|u|}/(2\omega)$; the case $P=p_1$ follows by applying it
to $-ih'$ and integrating by parts. Polynomial division by
$p_1^2+\omega^2$ gives the general case.

The transforms of the reflected and translated tests are
$-\theta_1\widehat b(-p_1,k)$ and
$e^{-ip_1t-ik\cdot y}\widehat b(p_1,k)$. For each pair of components,
the product of their time transforms is the transform of
\[
 h_{ac}(u;k)=\int
       \overline{\widehat b_{\perp,a}(s,k)}
       \widehat b_{\perp,c}(u-s-t,k)\,ds.
\]
This is smooth and compactly supported in $(0,\infty)$, because
$b$ is supported a positive distance from the mirror and $t\ge0$.
The full Fourier covariance is absolutely integrable, since $M$ is
bounded and the tests are Schwartz. Fubini permits integration first
in $p_1$. Applying the two preceding identities entrywise removes
$Q_k$ and replaces $N_k(p_1)$ by $N_k(-i\omega(k))$. Substitution
$u=s+s'+t$ factors the remaining integral into
$e^{-\omega(k)t}\overline{L_{b,a}(k)}L_{b,c}(k)$ and proves
\eqref{eq:example-time-kernel}.

At $t=0$, $y=0$ the integrand is nonnegative when $A\ge0$.
The Gaussian reflection criterion recalled in \cref{sec:examples}
therefore gives reflection positivity. The same specialization
shows that the displayed spectral measure is finite. The
reconstruction identity identifies the left side of
\eqref{eq:example-time-kernel} with the Fourier--Laplace transform
of the spectral measure of $[\B(b)]_1$. The right side is the
transform of the stated image measure, so uniqueness identifies them.
For completeness, uniqueness for finite complex measures follows by
first fixing integer $t\ge0$ and applying Fourier uniqueness in $y$.
After the change $s=e^{-\omega}$, equality of all these marginals
gives equality against every polynomial in $s$ on $[0,1]$.
Polynomial approximation gives equality against continuous functions;
rectangles then determine the original measures.
\end{proof}

\begin{lemma}[From one-field vectors to the Gaussian reconstruction]
\label{lem:example-gaussian-cone}
Let $\gamma$ be a centered Gaussian law invariant under translations
and reflection positive and invariant at $x_1=0$. Let
$K\subset[0,\infty)\times\R^2$ be closed, contain $(0,0)$ and
satisfy $K+K\subset K$. If the spectral measure of $[\B(b)]_1$
is carried by $K$ for every real smooth test compactly supported in
$\{x_1>0\}$, then the whole joint spectral measure $E_1$ is
carried by $K$.
\end{lemma}

\begin{proof}
Put $C(b,b')=\gamma(\B(b)\B(b'))$. For finitely many positive-side
real tests $b_a$ and coefficients $z_a$, let
\[
 \Phi=\sum_a z_a e^{i\B(b_a)},\qquad
 z'_a=z_a e^{-C(b_a,b_a)/2},\qquad \psi_a=[\B(b_a)]_1.
\]
Gaussian integration and the reconstruction identity give
\begin{align*}
 \langle[\Phi]_1,e^{-tH_1}V_1(y)[\Phi]_1\rangle
 &=\sum_{a,c}\overline{z'_a}z'_c
       \exp\bigl(C(\theta_1\cdot b_a,\tau_{te_1+y}\cdot b_c)\bigr),\\
 C(\theta_1\cdot b_a,\tau_{te_1+y}\cdot b_c)
 &=\int e^{-t\omega+ip\cdot y}\,dm_{ac}(\omega,p),
 \qquad m_{ac}=\langle\psi_a,E_1(\cdot)\psi_c\rangle.
\end{align*}
The measures $m_{ac}$ are finite and carried by $K$.
Indeed, Cauchy--Schwarz on finite Borel partitions gives
$|m_{ac}|(B)\le[m_{aa}(B)m_{cc}(B)]^{1/2}$.
On the additive semigroup $[0,\infty)\times\R^2$, the measure
\[
 \exp_*(m)=\sum_{r\ge0}\frac{m^{*r}}{r!},\qquad
 m^{*0}=\delta_{(0,0)},
\]
converges in total variation, since $\|m^{*r}\|\le\|m\|^r$.
Its Fourier--Laplace transform is the exponential of that of $m$.
The conditions on $K$ ensure that $\exp_*(m_{ac})$ is carried by
$K$. Transform uniqueness, as proved above, therefore identifies
the spectral measure of $[\Phi]_1$ with
\[
 \sum_{a,c}\overline{z'_a}z'_c\exp_*(m_{ac}).
\]
It follows that $E_1(K^c)[\Phi]_1=0$. The positive-side characters
span a dense subspace of the positive-side $L^2$ space, and their
classes are dense in the reconstruction; see
\cite[Section~7.1]{QEDCurrent}. Thus $E_1(K^c)=0$.
\end{proof}

\begin{thebibliography}{9}
\footnotesize
\bibitem{Dunkel} J.~Dunkel, \emph{18.04 Complex analysis with applications},
Spring 2019 lecture notes, adapted and extended from notes by
A.~Nachbin and J.~Orloff, Massachusetts Institute of Technology, 2019.
\href{https://math.mit.edu/~dunkel/Teach/18.04_2019S/notes/1804_Main.pdf}{Lecture notes}.
\bibitem{Melrose} R.~Melrose, \emph{Functional Analysis: Lecture notes
for 18.102, Spring 2020}, Massachusetts Institute of Technology, 2020.
\href{https://ocw.mit.edu/courses/18-102-introduction-to-functional-analysis-spring-2021/3d4cc88026d44a01f936cd6a0aa995cb_MIT18_102s20_lec_FA.pdf}{Lecture notes}.
\bibitem{QEDConsequences} \emph{Massless photons and Maxwell infrared limits
in three-dimensional QED}, companion manuscript, revised 30 October 2026.
\bibitem{Coincidences} \emph{Invariance at coinciding points}, companion
paper.
\bibitem{QEDCurrent} \emph{The field-strength Osterwalder--Schrader
axioms for QED in three dimensions}, comprehensive edition, 263~pp., 20~October 2026,
before the modular integration revision.

\end{thebibliography}
\end{document}
